\ifdefined\pdftexversion\pdfoutput=1\fi
% Options for packages loaded elsewhere
\PassOptionsToPackage{unicode}{hyperref}
\PassOptionsToPackage{hyphens}{url}
\PassOptionsToPackage{dvipsnames,svgnames,x11names}{xcolor}
\documentclass[
  10pt,
]{article}
\usepackage{xcolor}
\usepackage{amsmath,amssymb}
\setcounter{secnumdepth}{5}
\usepackage{iftex}
\ifPDFTeX
  \usepackage[T1]{fontenc}
  \usepackage[utf8]{inputenc}
  \usepackage{textcomp} % provide euro and other symbols
\else % if luatex or xetex
  \usepackage{unicode-math} % this also loads fontspec
  \defaultfontfeatures{Scale=MatchLowercase}
  \defaultfontfeatures[\rmfamily]{Ligatures=TeX,Scale=1}
\fi
\usepackage{lmodern}
\ifPDFTeX\else
  % xetex/luatex font selection
\fi
% Use upquote if available, for straight quotes in verbatim environments
\IfFileExists{upquote.sty}{\usepackage{upquote}}{}
\IfFileExists{microtype.sty}{% use microtype if available
  \usepackage[]{microtype}
  \UseMicrotypeSet[protrusion]{basicmath} % disable protrusion for tt fonts
}{}
\makeatletter
\@ifundefined{KOMAClassName}{% if non-KOMA class
  \IfFileExists{parskip.sty}{%
    \usepackage{parskip}
  }{% else
    \setlength{\parindent}{0pt}
    \setlength{\parskip}{6pt plus 2pt minus 1pt}}
}{% if KOMA class
  \KOMAoptions{parskip=half}}
\makeatother
\usepackage{color}
\usepackage{fancyvrb}
\newcommand{\VerbBar}{|}
\newcommand{\VERB}{\Verb[commandchars=\\\{\}]}
\DefineVerbatimEnvironment{Highlighting}{Verbatim}{commandchars=\\\{\}}
% Add ',fontsize=\small' for more characters per line
\newenvironment{Shaded}{}{}
\newcommand{\AlertTok}[1]{\textcolor[rgb]{1.00,0.00,0.00}{\textbf{#1}}}
\newcommand{\AnnotationTok}[1]{\textcolor[rgb]{0.38,0.63,0.69}{\textbf{\textit{#1}}}}
\newcommand{\AttributeTok}[1]{\textcolor[rgb]{0.49,0.56,0.16}{#1}}
\newcommand{\BaseNTok}[1]{\textcolor[rgb]{0.25,0.63,0.44}{#1}}
\newcommand{\BuiltInTok}[1]{\textcolor[rgb]{0.00,0.50,0.00}{#1}}
\newcommand{\CharTok}[1]{\textcolor[rgb]{0.25,0.44,0.63}{#1}}
\newcommand{\CommentTok}[1]{\textcolor[rgb]{0.38,0.63,0.69}{\textit{#1}}}
\newcommand{\CommentVarTok}[1]{\textcolor[rgb]{0.38,0.63,0.69}{\textbf{\textit{#1}}}}
\newcommand{\ConstantTok}[1]{\textcolor[rgb]{0.53,0.00,0.00}{#1}}
\newcommand{\ControlFlowTok}[1]{\textcolor[rgb]{0.00,0.44,0.13}{\textbf{#1}}}
\newcommand{\DataTypeTok}[1]{\textcolor[rgb]{0.56,0.13,0.00}{#1}}
\newcommand{\DecValTok}[1]{\textcolor[rgb]{0.25,0.63,0.44}{#1}}
\newcommand{\DocumentationTok}[1]{\textcolor[rgb]{0.73,0.13,0.13}{\textit{#1}}}
\newcommand{\ErrorTok}[1]{\textcolor[rgb]{1.00,0.00,0.00}{\textbf{#1}}}
\newcommand{\ExtensionTok}[1]{#1}
\newcommand{\FloatTok}[1]{\textcolor[rgb]{0.25,0.63,0.44}{#1}}
\newcommand{\FunctionTok}[1]{\textcolor[rgb]{0.02,0.16,0.49}{#1}}
\newcommand{\ImportTok}[1]{\textcolor[rgb]{0.00,0.50,0.00}{\textbf{#1}}}
\newcommand{\InformationTok}[1]{\textcolor[rgb]{0.38,0.63,0.69}{\textbf{\textit{#1}}}}
\newcommand{\KeywordTok}[1]{\textcolor[rgb]{0.00,0.44,0.13}{\textbf{#1}}}
\newcommand{\NormalTok}[1]{#1}
\newcommand{\OperatorTok}[1]{\textcolor[rgb]{0.40,0.40,0.40}{#1}}
\newcommand{\OtherTok}[1]{\textcolor[rgb]{0.00,0.44,0.13}{#1}}
\newcommand{\PreprocessorTok}[1]{\textcolor[rgb]{0.74,0.48,0.00}{#1}}
\newcommand{\RegionMarkerTok}[1]{#1}
\newcommand{\SpecialCharTok}[1]{\textcolor[rgb]{0.25,0.44,0.63}{#1}}
\newcommand{\SpecialStringTok}[1]{\textcolor[rgb]{0.73,0.40,0.53}{#1}}
\newcommand{\StringTok}[1]{\textcolor[rgb]{0.25,0.44,0.63}{#1}}
\newcommand{\VariableTok}[1]{\textcolor[rgb]{0.10,0.09,0.49}{#1}}
\newcommand{\VerbatimStringTok}[1]{\textcolor[rgb]{0.25,0.44,0.63}{#1}}
\newcommand{\WarningTok}[1]{\textcolor[rgb]{0.38,0.63,0.69}{\textbf{\textit{#1}}}}
\usepackage{longtable,booktabs,array}
\newcounter{none} % for unnumbered tables
\usepackage{calc} % for calculating minipage widths
% Correct order of tables after \paragraph or \subparagraph
\usepackage{etoolbox}
\makeatletter
\patchcmd\longtable{\par}{\if@noskipsec\mbox{}\fi\par}{}{}
\makeatother
% Allow footnotes in longtable head/foot
\IfFileExists{footnotehyper.sty}{\usepackage{footnotehyper}}{\usepackage{footnote}}
\makesavenoteenv{longtable}
\usepackage{graphicx}
\makeatletter
\newsavebox\pandoc@box
\newcommand*\pandocbounded[1]{% scales image to fit in text height/width
  \sbox\pandoc@box{#1}%
  \Gscale@div\@tempa{\textheight}{\dimexpr\ht\pandoc@box+\dp\pandoc@box\relax}%
  \Gscale@div\@tempb{\linewidth}{\wd\pandoc@box}%
  \ifdim\@tempb\p@<\@tempa\p@\let\@tempa\@tempb\fi% select the smaller of both
  \ifdim\@tempa\p@<\p@\scalebox{\@tempa}{\usebox\pandoc@box}%
  \else\usebox{\pandoc@box}%
  \fi%
}
% Set default figure placement to htbp
\def\fps@figure{htbp}
\makeatother
\setlength{\emergencystretch}{3em} % prevent overfull lines
\providecommand{\tightlist}{%
  \setlength{\itemsep}{0pt}\setlength{\parskip}{0pt}}
% TMLR journal style (tmlr.sty shipped alongside main.tex).
% Anonymous by default; switch to [accepted] for camera-ready or [preprint] to de-anonymize.
\usepackage{tmlr}
% arxiv.sty sets \rmdefault=ptm / \sfdefault=phv (Type1-era Times/Helvetica).
% Under tectonic's XeTeX path (TU encoding) those families don't exist, so every
% shape request — including italic — logs "font shape undefined" and substitutes
% the one default (LM Regular), silently flattening \emph. Give fontspec real
% OpenType clones with full shape sets instead: TeX Gyre Termes (Times) and
% Heros (Helvetica), both in the tectonic bundle.
% Latin Modern Mono (tectonic's \texttt) has no Greek/subscripts; tectonic warns
% and silently DROPS the glyph ("τ = 0.25" prints " = 0.25"). Map every math
% unicode char that can survive fold-math (inside code spans) to a math-mode
% equivalent so it renders instead of vanishing. Deliberately OUTSIDE the
% stripped XeTeX block: under arXiv's pdfLaTeX these chars are undefined and
% error, so the mapping serves both engines.
\usepackage{newunicodechar}
\newunicodechar{α}{\ensuremath{\alpha}}\newunicodechar{β}{\ensuremath{\beta}}
\newunicodechar{γ}{\ensuremath{\gamma}}\newunicodechar{δ}{\ensuremath{\delta}}
\newunicodechar{ε}{\ensuremath{\varepsilon}}\newunicodechar{ζ}{\ensuremath{\zeta}}
\newunicodechar{η}{\ensuremath{\eta}}\newunicodechar{θ}{\ensuremath{\theta}}
\newunicodechar{ι}{\ensuremath{\iota}}\newunicodechar{κ}{\ensuremath{\kappa}}
\newunicodechar{λ}{\ensuremath{\lambda}}\newunicodechar{μ}{\ensuremath{\mu}}
\newunicodechar{ν}{\ensuremath{\nu}}\newunicodechar{ξ}{\ensuremath{\xi}}
\newunicodechar{π}{\ensuremath{\pi}}\newunicodechar{ρ}{\ensuremath{\rho}}
\newunicodechar{σ}{\ensuremath{\sigma}}\newunicodechar{τ}{\ensuremath{\tau}}
\newunicodechar{φ}{\ensuremath{\varphi}}\newunicodechar{χ}{\ensuremath{\chi}}
\newunicodechar{ψ}{\ensuremath{\psi}}\newunicodechar{ω}{\ensuremath{\omega}}
\newunicodechar{Δ}{\ensuremath{\Delta}}\newunicodechar{Σ}{\ensuremath{\Sigma}}
\newunicodechar{Ω}{\ensuremath{\Omega}}\newunicodechar{Γ}{\ensuremath{\Gamma}}
\newunicodechar{₀}{\ensuremath{_0}}\newunicodechar{₁}{\ensuremath{_1}}
\newunicodechar{₂}{\ensuremath{_2}}\newunicodechar{₃}{\ensuremath{_3}}
\newunicodechar{₄}{\ensuremath{_4}}\newunicodechar{₅}{\ensuremath{_5}}
\newunicodechar{₆}{\ensuremath{_6}}\newunicodechar{₇}{\ensuremath{_7}}
\newunicodechar{₈}{\ensuremath{_8}}\newunicodechar{₉}{\ensuremath{_9}}
\newunicodechar{ᵢ}{\ensuremath{_i}}\newunicodechar{ᵥ}{\ensuremath{_v}}
\newunicodechar{⁰}{\ensuremath{^0}}\newunicodechar{¹}{\ensuremath{^1}}
\newunicodechar{²}{\ensuremath{^2}}\newunicodechar{³}{\ensuremath{^3}}
\newunicodechar{⁻}{\ensuremath{^-}}\newunicodechar{⁺}{\ensuremath{^+}}
\newunicodechar{≈}{\ensuremath{\approx}}\newunicodechar{≠}{\ensuremath{\neq}}
\newunicodechar{∝}{\ensuremath{\propto}}\newunicodechar{∞}{\ensuremath{\infty}}
\newunicodechar{±}{\ensuremath{\pm}}\newunicodechar{‖}{\ensuremath{\|}}
\newunicodechar{−}{\ensuremath{-}}\newunicodechar{⋅}{\ensuremath{\cdot}}
\usepackage{etoolbox}
% fold-math.py emits \argmin/\argmax as named operators; amsmath doesn't define them.
\DeclareMathOperator*{\argmin}{arg\,min}
\DeclareMathOperator*{\argmax}{arg\,max}
\AtBeginEnvironment{longtable}{\footnotesize\setlength{\tabcolsep}{4pt}}
\setlength{\emergencystretch}{3em}

% Capability-table glyphs. Checkmarks/crosses via pifont (\ding, a matched pair);
% Harvey balls via tikz (clash-free with amssymb, which wasysym is not). Each is
% wrapped in accsupp ActualText so it extracts as its Unicode char (pdftotext,
% screen readers, copy-paste) instead of a dingbat-font artifact.
\usepackage{pifont}
\usepackage{tikz}
\usepackage{accsupp}
\newcommand{\axglyph}[2]{\BeginAccSupp{method=hex,unicode,ActualText=#1}#2\EndAccSupp}
\DeclareRobustCommand{\cmark}{\axglyph{2713}{\ding{51}}}
\DeclareRobustCommand{\xmark}{\axglyph{2717}{\ding{55}}}
\DeclareRobustCommand{\hvfull}{\axglyph{25CF}{\tikz[baseline=-0.5ex]\fill (0,0) circle (0.45ex);}}
\DeclareRobustCommand{\hvempty}{\axglyph{25CB}{\tikz[baseline=-0.5ex]\draw[line width=0.35pt] (0,0) circle (0.45ex);}}
\DeclareRobustCommand{\hvhalf}{\axglyph{25D0}{\tikz[baseline=-0.5ex]{\fill (0,0.45ex) arc (90:270:0.45ex) -- cycle; \draw[line width=0.35pt] (0,0) circle (0.45ex);}}}
\usepackage{bookmark}
\IfFileExists{xurl.sty}{\usepackage{xurl}}{} % add URL line breaks if available
\urlstyle{same}
% fallback for those not using the hyperref driver hyperxmp:
\makeatletter
\@ifundefined{xmpquote}{\newcommand{\xmpquote}[1]{#1}}{}
\makeatother
\hypersetup{
  colorlinks=true,
  linkcolor={blue},
  filecolor={Maroon},
  citecolor={Blue},
  urlcolor={blue},
  pdfcreator={LaTeX via pandoc}}

\title{The Hypothesis Graph: A Verifiable Semantic Memory for Coding Agents}
% TMLR author block: hidden while tmlr.sty runs without [accepted]/[preprint].
\author{\name June Kim \email june@june.kim \\
      \addr Independent Researcher}

\begin{document}
\maketitle
\begin{abstract}
The \textbf{hypothesis graph} is a data structure for coding agents that
deepens their reasoning and makes them accountable. Implemented at the
harness layer, its nodes are testable claims, its edges the refutations
that name the next claim. It updates by inquiry, Peirce's typed loop of
abduction, deduction, and induction. What is new is running that loop as
mechanical harness operations around the one step no procedure reaches.
That step, the abductive leap, stays in the model; the harness
manufactures the surprise, fires the kill, and records the trail,
instead of leaving the loop to a model's undifferentiated prose. We
demonstrate the mechanism on one contamination-free bugfix. Six
self-attested arms plateau at a narrow fix; an externally verified
comparator, the one input a context-bound agent cannot author for
themselves, carries a weaker model (Sonnet 4.6) to the behavior of the
merged human fix that the strongest released model (Fable) does not
reach without it. This buys verifiable accountability and reasoning that
persists past the context window, at no additional training cost, usable
by any coding-agent harness that can run pinned trials and read their
verdicts.
\end{abstract}

\section{Introduction}\label{introduction}

In coding agents, the LLM is wrapped in a harness: the verification,
testing, and memory a software task needs. The field building these is
moving up a level of abstraction, from the model to the harness.
Roychoudhury et al.~(2025) reframe the goal as \emph{programming with
trust}, arguing that deployment turns on verification, testing, and
analysis built into the agent rather than on raw generation
(\href{https://arxiv.org/abs/2502.13767}{arXiv:2502.13767}). Liu et
al.~(2024) survey agents for software engineering and organize the field
around those same missing pieces
(\href{https://arxiv.org/abs/2409.02977}{arXiv:2409.02977}); Yehudai et
al.~(2025) add that scoring final outputs misses the reasoning and
failure causes inside a run, and call for trajectory-level assessment
(\href{https://arxiv.org/abs/2503.16416}{arXiv:2503.16416}); Wang et
al.~(2025) survey agentic-programming systems and list persistent,
structured memory among the open challenges
(\href{https://arxiv.org/abs/2508.11126}{arXiv:2508.11126}).

A patch passes the visible tests, but that's not enough: passing
certifies only the cases the tests cover. An over-narrow patch passes
them and is wrong off-suite. Confirming it means reconstructing the
reasoning the agent never recorded, at a cost approaching that of
producing it, so the work shifts from writing to checking. Code review
is the bottleneck.

When reasoning is discarded, each run rebuilds context from scratch.
Even where agent memory adopts the cognitive-architecture lineage, as
CoALA (Sumers et al.~2024) does in mapping it onto Soar (Laird 1987) and
ACT-R, the semantic slot stores facts rather than a falsifiable
structure, so the search is discarded once a patch passes. At best, a
trail of blobs is saved as provenance.

We were promised a junior developer with near-infinite patience. All we
got was tool calls in a loop: a cracked-up amnesiac contractor, leaving
mistakes for maintainers to review.

{\def\LTcaptype{none} % do not increment counter
\begin{longtable}[]{@{}
  >{\raggedright\arraybackslash}p{(\linewidth - 2\tabcolsep) * \real{0.2500}}
  >{\raggedright\arraybackslash}p{(\linewidth - 2\tabcolsep) * \real{0.7500}}@{}}
\toprule\noalign{}
\begin{minipage}[b]{\linewidth}\raggedright
Problem
\end{minipage} & \begin{minipage}[b]{\linewidth}\raggedright
Solution
\end{minipage} \\
\midrule\noalign{}
\endhead
\bottomrule\noalign{}
\endlastfoot
\emph{Reasoning cost} & A domain-aware \emph{division of intellect}: the
model reasons, the harness holds the open hypotheses and runs the
gate. \\
\emph{Agent trust} & A \emph{verifiable reasoning chain}: a claim clears
only when an independent party reruns its recorded trial. \\
\emph{Lost reasoning} & A \emph{human-legible representation}: typed
nodes a later run or a human reads and replays. \\
\emph{Review bottleneck} & Every step \emph{grounded in a trial}: a
reviewer accepts the work or resumes where it stuck. \\
\end{longtable}
}

Here we give the agent a trail of \textbf{verifiable knowledge}. The
\textbf{hypothesis graph} is a semantic memory that holds an agent's
reasoning while it works. The fix arrives with the inquiry that produced
it: what was hypothesized, what was tested, what was ruled out. Each
step carries a trial a stranger can rerun. The promise reverses the
trust assumption by shifting the burden of proof outside the agent. The
model still reasons, but each node it produces is checkable on its own
at the harness layer, independent of the model.

Two contributions run through this, and they are separable. The graph is
the accountability substrate, and its value is argued by construction:
every node replays. The externalized comparator is the capability
mechanism, and its value is shown by ablation. The graph makes reasoning
checkable; the comparator supplies the one verdict a self-graded agent
cannot author for itself. Keeping them apart is what lets the ablation
below read cleanly, where several arms share the graph and only the
verdict source moves.

What that buys is not incremental. On one contamination-free bug, an
externalized comparator carried a weaker model to a fix the strongest
released model could not reach without it (\S\ref{right-regime}): a
capability lift, where a scaffold usually buys only automation.

\section{The hypothesis graph}\label{hygraph}

\subsection{Requirements}\label{requirements}

A semantic memory that holds an agent's live reasoning must clear four
requirements at once:

\begin{itemize}
\tightlist
\item
  \textbf{Holds hypotheses.} Stores a hypothesis still under
  consideration, where prose, skill libraries, and retrieval keep only
  verified facts and established chunks: the hypothesis-shaped gap none
  of them fills.
\item
  \textbf{Refutable by test.} Each claim carries a kill condition: the
  executable test that can prove it wrong.
\item
  \textbf{Independently verifiable.} A stranger verifies a conclusion by
  rerunning its recorded trial, instead of re-deriving the reasoning or
  taking the author's word.
\item
  \textbf{Persistent memory.} The trail persists past the context window
  rather than being discarded once a patch passes.
\end{itemize}

Prior structures fill some of these roles (\cmark{} native, \hvhalf{}
with common extensions, \xmark{} needs another layer doing the work).

{\def\LTcaptype{none} % do not increment counter
\begin{longtable}[]{@{}
  >{\raggedright\arraybackslash}p{(\linewidth - 8\tabcolsep) * \real{0.4306}}
  >{\centering\arraybackslash}p{(\linewidth - 8\tabcolsep) * \real{0.1250}}
  >{\centering\arraybackslash}p{(\linewidth - 8\tabcolsep) * \real{0.1111}}
  >{\centering\arraybackslash}p{(\linewidth - 8\tabcolsep) * \real{0.1944}}
  >{\centering\arraybackslash}p{(\linewidth - 8\tabcolsep) * \real{0.1389}}@{}}
\toprule\noalign{}
\begin{minipage}[b]{\linewidth}\raggedright
Structure
\end{minipage} & \begin{minipage}[b]{\linewidth}\centering
Holds hypotheses
\end{minipage} & \begin{minipage}[b]{\linewidth}\centering
Includes tests
\end{minipage} & \begin{minipage}[b]{\linewidth}\centering
Independently verifiable
\end{minipage} & \begin{minipage}[b]{\linewidth}\centering
Persistent memory
\end{minipage} \\
\midrule\noalign{}
\endhead
\bottomrule\noalign{}
\endlastfoot
\textbf{Hypothesis graph} (this work) & \cmark{} & \cmark{} & \cmark{} &
\cmark{} \\
Truth-maintenance (Doyle 1979; de Kleer 1986) & \cmark{} & \hvhalf{} &
\hvhalf{} & \hvhalf{} \\
Provenance / lineage (W3C PROV, Moreau et~al.~2013) & \hvhalf{} &
\hvhalf{} & \hvhalf{} & \cmark{} \\
Search + proof tree (Clarke et~al.~2000; Solar-Lezama et~al.~2006) &
\cmark{} & \hvhalf{} & \hvhalf{} & \hvhalf{} \\
Argumentation (Dung 1995; Modgil \& Prakken 2014) & \cmark{} & \hvhalf{}
& \hvhalf{} & \hvhalf{} \\
Event-sourced log / ReAct trace (Yao et~al.~2023) & \cmark{} & \hvhalf{}
& \hvhalf{} & \cmark{} \\
\end{longtable}
}

\emph{The middle two columns are verification (a falsifiable trial a
stranger can rerun), the outer two are the hypothesis-shaped gap and
retention. Only the hypothesis graph clears all four in one append-only
file.}

An LLM is what finally fills it. Filling the graph takes a reasoner that
reads a surprising failure, proposes candidate causes in open
vocabulary, and turns each into an executable test, with no hand-built
domain model. Classical inference engines could do this only inside a
formalism encoded by hand, which is why the slot stayed a research
program; the LLM populates it across arbitrary codebases, which is what
makes the structure practical here.

\subsection{Graph semantics}\label{graph-semantics}

The structure meets all four with three pieces, a node, an edge, and one
invariant. A \textbf{node} is a claim bound to a trial: a hypothesis,
the perturbation that tests it stated as an exact command, the observed
outcome, and a credence typed by the reasoning mode that established it.
An \textbf{edge} is generated by a kill condition: the manner of a
hypothesis's death names the next hypothesis, so the structure is
self-extending, with no external controller deciding where to look. The
invariant comes before the operations, because every operation is
defined to preserve it:

\begin{quote}
\textbf{Replay invariant.} Every node is reconstructible, by a stranger
who does not trust the author, from its recorded trial alone: the exact
command, the observed outcome, the verdict, and the credence cap.
\end{quote}

The contract holds on the node's mechanical skeleton; the hypothesis
prose the node also carries falls outside it, by design. The prose is
the part an auditor would otherwise have to trust, and the recorded
trial is what replaces trusting it, so the invariant draws its line
exactly where checkability begins. The guarantee is narrow and named:
the command, outcome, and verdict are checkable, while the mode label
that caps a credence is a convention the writer is trusted to apply
honestly.

\subsection{Operations}\label{operations}

Five operations maintain the structure, each defined with the one-clause
argument that it preserves the invariant, the way a balanced tree's
insert is defined to restore balance:

\begin{itemize}
\tightlist
\item
  \textbf{Create} (a smart constructor), append a node: abduction writes
  a hypothesis with its kill condition and the exact trial that tests it
  (``the bulb is dead'', trial \texttt{swap\ in\ a\ fresh\ bulb}).
  \emph{Preservation}: a node is admitted only if its trial replays, so
  by construction every committed node satisfies the contract, and the
  invariant holds by induction over reachable graphs. This admission
  rule is what makes every later operation safe.
\item
  \textbf{Read / replay}: ask which hypotheses are still open, and
  reconstruct any node by re-running its recorded trial, no trust in the
  author required. \emph{Preservation} is vacuous, read mutates nothing,
  but replay is the one that does the work, the check every other
  operation's preservation is stated against.
\item
  \textbf{Classify} (the update): a trial's outcome marks its node
  \emph{killed} or \emph{witnessed} and caps its credence at the mode
  that earned it, a verdict written once. \emph{Preservation}: classify
  appends a verdict and never edits the recorded trial, so the node
  still replays to the same outcome.
\item
  \textbf{Link} (edge-from-kill): the manner of a hypothesis's death
  names the next hypothesis (the dead bulb names ``the dimmer is
  broken'', trial \texttt{bypass\ the\ dimmer\ to\ the\ wall}), so a
  classification spawns the next node with no external controller.
  \emph{Preservation}: link only appends, and each new node is admitted
  under Create's rule. This is the one genuinely novel operation.
\item
  \textbf{Prune}: a dead branch leaves the working frontier while its
  record stays in place. \emph{Preservation} is trivial, prune is a
  frontier-set operation that changes what is \emph{live} and deletes
  nothing, so every pruned node replays exactly as before.
\end{itemize}

\begin{figure}
\centering
\pandocbounded{\includegraphics[keepaspectratio,alt={A hypothesis graph, two nodes and the edge between them, on the dead-light inquiry. The bulb hypothesis is killed by a cheap trial (swap in a fresh bulb, still dark); its death names the next node, the dimmer, which a second trial witnesses (bypass it to the wall, the light comes on). Each node binds a hypothesis to a trial, an observed outcome, and a credence capped by the mode that earned it: abduction proposes and stays low, induction is test-backed and rises. Every node rebuilds from its recorded trial, so an auditor replays the structure instead of trusting it.}]{hypothesis-graph-anatomy.pdf}}
\caption{A hypothesis graph, two nodes and the edge between them, on the
dead-light inquiry. The bulb hypothesis is killed by a cheap trial (swap
in a fresh bulb, still dark); its death names the next node, the dimmer,
which a second trial witnesses (bypass it to the wall, the light comes
on). Each node binds a hypothesis to a trial, an observed outcome, and a
credence capped by the mode that earned it: abduction proposes and stays
low, induction is test-backed and rises. Every node rebuilds from its
recorded trial, so an auditor replays the structure instead of trusting
it.}
\end{figure}

\subsection{Auditability}\label{auditability}

The invariant yields an essential property that no benchmark can
evaluate:

\begin{quote}
\textbf{Local Replay Auditability.} Any single conclusion is checkable
by re-running that node's one recorded trial, without reconstructing the
inquiry and without extending trust to whoever ran it.
\end{quote}

This is, for inquiry, the analogue of a certificate a consumer checks
without trusting the producer, and like the Merkle audit path or
proof-carrying code, its value rests on a contract; a complexity bound
is beside the point.

Two grades of it matter. Where the trial is a deterministic command over
pinned inputs, a compiler on a fixed toolchain, a test in a container,
replay is \emph{strong}: re-execution reproduces the recorded outcome,
and the lead case (\S\ref{right-regime}) is here. Where the trial runs a
model or a live service, replay is \emph{artifact-level}: the recorded
output is verified and the deterministic predicate re-run over it, the
command standing as a provenance event rather than a reproducible
computation. The honest scope is the first shading into the second as
the system under test grows less deterministic. Pruning leaves all of it
untouched: a pruned branch drops from the working frontier but stays in
the record, so its nodes replay exactly as before.

\subsection{Knowledge maintenance}\label{knowledge-maintenance}

The nodes are ordinary; what is novel is the edge semantics. A search
tree finds; a proof tree justifies. The hypothesis graph is both at
once, because the search path \emph{is} the justification: every step
was a trial.

It sits at the confluence of older lineages: truth-maintenance and
model-based diagnosis (de Kleer 1986; Reiter 1987; de Kleer \& Williams
1987), sequential experimental design (Wald 1947; Vovk \& Wang 2021),
abstract argumentation (Dung 1995), and counterexample-guided refinement
(CEGAR, Clarke et al.~2000; CEGIS, Solar-Lezama et al.~2006). Refinement
is the closest kin: a counterexample \emph{is} a kill that names the
next experiment, the hypothesis graph's defining edge. What is novel is
running it over an \emph{open} hypothesis space abduced in domain
vocabulary. Replayability stands in for the sound abstraction a closed
setting supplies for free, with completeness as the price the open move
forfeits (\S\ref{lineage}).

The near neighbors each hold part of this and source the rest from
outside themselves. A truth-maintenance system (Doyle 1979; de Kleer
1986) maintains belief status under assumptions, but the empirical trial
and the kill-generated successor are external conventions. Provenance
(W3C PROV, Moreau et al.~2013) records replayable activities, yet does
not decide which hypothesis comes next. And the ReAct trace (Yao et
al.~2023), the strongest mundane baseline, is an append-only log whose
continuation policy the controller decides and the record never holds.

What the hypothesis graph adds is their composition in one append-only
object: an open-domain hypothesis, its executable trial, its kill
condition, and the successor that kill names. That object belongs to the
verifiable family whose value is a contract rather than a complexity
bound, certificate transparency (Laurie et al., RFC 9162),
proof-carrying code (Necula 1997), content-addressed provenance: a data
structure paired with the protocol that writes and checks it.

\subsection{Semantic memory}\label{semantic-memory}

This is the data structure for \emph{testable} inquiry, and its entire
power is the perturbation surface. Strip the ability to poke the system
and read an outcome, and the same shape degrades into a plausibility
tree, which is the confabulation failure mode it exists to prevent.
Intuition is not verifiable from outside, so inquiry that has to be
checked trades it for an explicit perturbation surface. The hypothesis
graph is the verifiable serialization reasoning compiles to, so it can
be checked by someone who does not trust you. Proof is to intuition as
the hypothesis graph is to inquiry: not the thinking, the residue of the
thinking that survives a stranger's replay.

That residue is what the memory typology calls the \texttt{smem}:
persistent, typed, queryable, and owned by the harness rather than the
model. Those same properties make it an interface for agent interop:
because the structure is typed and external, a second agent, a later
run, or a human auditor reads and writes against one contract and can
rerun any node rather than trust it (\S\ref{discussion}). The graph in
this work is one markdown file per inquiry.

\section{Actionable epistemology}\label{epistemics}

Knowing is an act to update the credence of a claim. This is the subject
of \emph{Verifiable Knowledge}; the property that matters here, a
verdict a model cannot author for itself, is the one
\S\ref{right-regime} turns into a capability lift. In the Hypothesis
Graph, knowledge has the following properties:

\begin{itemize}
\tightlist
\item
  \textbf{Three states.} Witnessed is true, a build presently passing;
  killed is false, a build gone red; open is \emph{untrue}, a conjecture
  awaiting its test. These are the states of the entitlement ledger
  developed in \emph{Verifiable Knowledge}, scoped here to one node and
  its kill edge.
\item
  \textbf{Credence.} A node carries a credence capped by the mode that
  earned it, low for abduction, higher once tested (Ramsey 1926), the
  step a bare LLM skips when it emits uniform confidence with no
  propagation along the chain.
\item
  \textbf{Survived belief.} A node counts as knowledge only after
  withstanding a trial, a verificationist criterion (Ayer 1936), indexed
  to the stakes of acting on it. Before trial, a claim is a hypothesis.
\item
  \textbf{Causal connections.} Each node wires to what it depends on, so
  claims compose into a derivation, every step caused by an earlier one
  (Pearl 1988).
\end{itemize}

With each round of inquiry the causal boundary sharpens, and a fix that
respects it generalizes past the case at hand. With each successive
trial, credence accumulates to actionable confidence.

\begin{figure}
\centering
\pandocbounded{\includegraphics[keepaspectratio,alt={A claim and its trial as one record. Left, the claim: a load resting on a span, this bridge holds the load. Right, the trial: a toothpick-and-gumdrop model bridge bearing a steel weight across two supports. Photo: Oregon Department of Transportation, CC BY 2.0.}]{bridge-trial-light.pdf}}
\caption{A claim and its trial as one record. Left, the claim: a load
resting on a span, \emph{this bridge holds the load}. Right, the trial:
a toothpick-and-gumdrop model bridge bearing a steel weight across two
supports. Photo: Oregon Department of Transportation, CC BY 2.0.}
\end{figure}

This is one more projection of the protocol \emph{Verifiable Knowledge}
sets out, from where the agent stands, and its payoff is transfer.
Paired with its trial, a node carries its warrant across intact; handed
on the author's word, only the verdict crosses and the warrant is
re-derivable from scratch.

A coding agent is already a strong internal verifier, running the test
and verifying that it passed. What it lacked was the protocol, extending
credence beyond the context window. With the protocol, knowledge
compounds for other agents and humans.

\section{Methodeutics: a discipline of inquiry}\label{grounding}

A store of causal knowledge needs reasoning to be encoded into it. How
can reasoning become mechanical enough for encoding? The insight, the
leap to a candidate cause, stays with the model, but the harness needs a
method to stage reasoning. Reasoning then becomes mechanical the way a
proof is, through the discipline of checking ideas. In a precise and
limited sense the harness stages and checks reasoning mechanically,
everything but the leap, which it can only trigger. This discipline,
Peirce called \textbf{inquiry}.

His \emph{Illustrations of the Logic of Science} (1878) and
\emph{Pragmatism as the Logic of Abduction} (1903) type the operations
of inquiry into three irreducible modes.

\begin{itemize}
\tightlist
\item
  \textbf{Abduction} generates explanatory hypotheses for surprising
  observations: \emph{what would, if true, make this no longer
  surprising?}
\item
  \textbf{Deduction} derives testable predictions from hypotheses:
  \emph{if this hypothesis holds, what follows?}
\item
  \textbf{Induction} tests predictions against evidence: \emph{does the
  evidence accord with the prediction?}
\end{itemize}

\begin{figure}
\centering
\pandocbounded{\includegraphics[keepaspectratio,alt={The three modes as one cycle: Observation \textbackslash rightarrow Theory (abduction), Theory \textbackslash rightarrow Experiment (deduction), Experiment \textbackslash rightarrow Observation (induction).}]{modes-of-reason-triangle-light.pdf}}
\caption{The three modes as one cycle: Observation \(\rightarrow\)
Theory (abduction), Theory \(\rightarrow\) Experiment (deduction),
Experiment \(\rightarrow\) Observation (induction).}
\end{figure}

No single mode carries a belief to its grade. Abduction proposes content
but does not test it; induction tests but introduces no new explanatory
content; deduction traces consequences but invents nothing. The credence
a node ends up with is what traversing all three earns it, and that is
what it means to call the modes typed: each is fixed by what it can't
do.

This is where readers balk: why is \emph{abduction} responsible for
theory? Isn't that deduction's job? No.~Deduction neither generates the
theory nor proves it; it unfolds the hypothesis into the predictions it
must answer for. The theory was abduced, the predictions deduced, and
induction does the testing.

Keep them separate and each does its one job; collapse them and you get
familiar failure modes:

\begin{itemize}
\tightlist
\item
  \textbf{Confirmation bias}: induction without abductive alternatives
\item
  \textbf{Confabulation}: abduction without inductive grounding
\item
  \textbf{Free-association}: no typed mode at all
\end{itemize}

That collapse is exactly what modern LLM agents do by default, since a
single forward pass proposes, predicts, evaluates, and rationalizes in
undifferentiated prose. Methodeutics, Peirce's term for the methodology
of inquiry, is how to conduct the typed-mode loop well. Encoded as
skills, it constructs and maintains the \texttt{smem}.

\emph{Modes of reason and the irreducible three.} Around the act of
testing, philosophy of science built an apparatus of real rigor: Bacon's
induction (1620), Popper's falsifiability (1934), Meehl's ``soft
science'' critique (1967), Pearl's causal calculus (2009). Justification
got its method, every step of it. But it begins one step too late,
taking the hypothesis as given and filing its origin under inspiration.
The discipline built an epistemology of justification and little of
discovery. Peirce named the missing operation, abduction; the discipline
still filed the origin of hypotheses under inspiration. The harness
gives it a first-class typed slot and triggers it, though the leap
itself stays the model's.

\section{Methodeutics, applied}\label{application}

Putting the theory to work means generating hypotheses as typed nodes
the harness can test, instead of trusting whatever a model guesses.

\emph{The surprise has a primitive, and it is a diff.} What the harness
manufactures is not the hypothesis but the discrepancy that provokes
one: a before snapshot, an after snapshot, and the perturbation that
flipped, read as figure against the ground that held (Rubin's Gestalt
terms). Separation logic calls the frame-inference half of this
\emph{bi-abduction} and scaled it to real codebases in Facebook Infer
(Calcagno et al.~2009; O'Hearn 2019), borrowing Peirce's word for an
operation that is not his abduction: it localizes where belief and code
diverge, the surprise, and leaves the leap to what would explain it
unmade. \texttt{inquire} works at the simplest level: one before/after
diff, the frame inferred from the symptom. The extensions to branches
and compositional cases are in the \hyperref[lineage]{lineage appendix}.

The ``XOR'' used throughout is shorthand for that separation: the figure
(what the fix must change) held apart from the ground (the frame that
stays invariant). Where bi-abduction infers the frame to make a proof go
through, the harness fires the same split as a check, computing the
symmetric difference against a known-good oracle and keeping only the
cases it flags. The XOR is the surprise, not the leap: it marks what a
hypothesis must explain, and the explaining stays the model's.

\begin{figure}
\centering
\pandocbounded{\includegraphics[keepaspectratio,alt={Bi-abduction on a dead fixture. With dimmer, fixture, and bulb all intact, the static scene names no suspect; the perturbation bypasses the dimmer to the wall, and the XOR isolates the figure (the dimmer) from the ground (fixture and bulb).}]{bi-abduction-dimmer.pdf}}
\caption{Bi-abduction on a dead fixture. With dimmer, fixture, and bulb
all intact, the static scene names no suspect; the perturbation bypasses
the dimmer to the wall, and the XOR isolates the figure (the dimmer)
from the ground (fixture and bulb).}
\end{figure}

\emph{Directed graphs as reasoning representation.} Pearl 1988
(\emph{Probabilistic Reasoning in Intelligent Systems}; Bayesian
networks as DAGs of dependencies); Pearl 2000/2009 (\emph{Causality};
structural causal models, d-separation, do-calculus). Pearl's lineage
was built for causal-structure inference; our data structure (typed
nodes, directed edges) puts it to hypothesis representation. The
difference from a Bayesian network is one of kind, and runs deeper than
dropped probabilities. A Bayes net conditions over a fixed variable set
and propagates probability along edges of dependence; the hypothesis
graph abduces its nodes as the inquiry runs, its edges genealogical, a
dead hypothesis naming its successor rather than a conditional
dependence. A Bayes net is justification over a space it is handed; the
hypothesis graph generates the space.

\begin{figure}
\centering
\pandocbounded{\includegraphics[keepaspectratio,alt={The hypothesis graph for the dead fixture. Abduction fans the observation into four typed candidate nodes; mechanical kill predicates fire on three (the socket, fixture, and bulb each cleared by a cheap test), the dimmer node is witnessed by the bypass and closes the last open hypothesis, and deduction derives the fix. Typed nodes, directed edges, all three modes in one inquiry.}]{hypothesis-graph-fixture.pdf}}
\caption{The hypothesis graph for the dead fixture. Abduction fans the
observation into four typed candidate nodes; mechanical kill predicates
fire on three (the socket, fixture, and bulb each cleared by a cheap
test), the dimmer node is witnessed by the bypass and closes the last
open hypothesis, and deduction derives the fix. Typed nodes, directed
edges, all three modes in one inquiry.}
\end{figure}

Isn't generating that space just debugging? It is, and debugging has
been automated for decades: spectrum-based fault localization,
statistical and delta debugging, model-based diagnosis, and search-based
program repair are mature fields (Jones et al.~2002; Liblit et al.~2005;
Zeller \& Hildebrandt 2002; Reiter 1987; Le Goues et al.~2012; Monperrus
2018). Debugging tools already automate the loop; what is new is that
the hypothesis graph persists it. Every engineer, and every repair tool,
runs some version of abduce a cause, kill it on evidence, witness the
survivor, derive the fix. But the engineer runs it in their head, and
the tools, whatever logs they keep, do not persist the search as a
typed, replayable hypothesis graph.

We implement this loop as a tool. \emph{comparator}
(\href{https://anonymous.4open.science/r/comparator-DA10}{anonymous.4open.science/r/comparator-DA10})
externalizes the surprise, the diff generation, outside the context
window, so a model has to represent the rule instead of tabulating the
case in front of it: it enumerates a space wider than the model's
hypothesis, calibrates each case against a known-good baseline, and
exposes one pass/fail gate, with the answer key held outside the model's
view. A failing case is a counterexample that forces the next fix, the
model's own leap, and the search records itself as the hypothesis graph,
fixes as nodes and counterexamples as edges.

\section{The methodeutic harness}\label{method}

How do epistemology and debugging become an agentic harness? The loop
that writes the graph needs four things:

\begin{itemize}
\tightlist
\item
  \textbf{Iteration.} One pass can't be trusted, so the loop re-enters
  on failure.
\item
  \textbf{Deterministic oracle.} A model can't be trusted to catch its
  own bias, so the check comes from outside the weights.
\item
  \textbf{Three modes of reasoning.} Abduction, deduction, and induction
  stay typed and separate, each capped at its own confidence.
\item
  \textbf{Semantic memory.} The reasoning is recorded into the
  hypothesis graph and survives the context window.
\end{itemize}

Concretely, this is a skill with a tool call in a loop: the outer
deterministic driver invokes an agent via the \texttt{inquire} skill,
who accesses the deterministic \texttt{comparator} tool. It mechanizes the
surprise, the diff the model would otherwise have to compute by hand.

\begin{figure}
\centering
\pandocbounded{\includegraphics[keepaspectratio,alt={The inquire skill: Peirce's three modes as a procedure that writes the hypothesis graph. Induction fires a deterministic kill or witness with no model arbitrating. implement and attest, which read the survivors and verify the patch, follow below.}]{inquire-skill.pdf}}
\caption{The \protect\texttt{inquire} skill: Peirce's three modes as a
procedure that writes the hypothesis graph. Induction fires a
deterministic kill or witness with no model arbitrating.
\protect\texttt{implement} and \protect\texttt{attest}, which read the
survivors and verify the patch, follow below.}
\end{figure}

\subsection{The inquiry frame}\label{inquiry-frame}

We recast each issue as an inquiry on an engineered system: a failure
trace, a codebase, a root cause to find, and an intervention that must
not regress the rest of the system. Code is the right substrate for the
hypothesis graph because it combines three properties that other inquiry
domains rarely bring together:

\begin{itemize}
\tightlist
\item
  \textbf{Reproducible}: same input yields same output, modulo
  controlled nondeterminism
\item
  \textbf{Deterministic}: causal lines from input to behavior are
  mechanical
\item
  \textbf{Perturbable}: single-line and single-function diffs are cheap
  to apply and fully observable
\end{itemize}

Because those three hold together, kill conditions over code are exact
executions.

One trial settles the predicate in this regime. In code the per-case
response is mechanically observable, so a single passing test on a
captured diff is a complete verdict that the diff satisfies the
executable predicate for that case. That verdict speaks to the predicate
alone: behaviors it doesn't cover are out of scope, a boundary that
\S\ref{right-regime} shows is exactly where the differences live. Where
such verdicts are aggregated, the right summary is counts and
denominators rather than confidence intervals: per-case verdicts are
exact, and aggregating them is bookkeeping.

The three Peircean modes are how \texttt{inquire} builds the graph, each
node typed by the mode that established it and capped at that mode's
confidence:

{\def\LTcaptype{none} % do not increment counter
\begin{longtable}[]{@{}
  >{\raggedright\arraybackslash}p{(\linewidth - 4\tabcolsep) * \real{0.1667}}
  >{\raggedright\arraybackslash}p{(\linewidth - 4\tabcolsep) * \real{0.7083}}
  >{\raggedright\arraybackslash}p{(\linewidth - 4\tabcolsep) * \real{0.1250}}@{}}
\toprule\noalign{}
\begin{minipage}[b]{\linewidth}\raggedright
Mode
\end{minipage} & \begin{minipage}[b]{\linewidth}\raggedright
What \texttt{inquire} does
\end{minipage} & \begin{minipage}[b]{\linewidth}\raggedright
Confidence
\end{minipage} \\
\midrule\noalign{}
\endhead
\bottomrule\noalign{}
\endlastfoot
\textbf{Abduction} & Proposes candidate root causes from the observed
failure; writes hypothesis nodes with falsifiable predicates and kill
conditions (read-only) & low \\
\textbf{Deduction} & Traces each hypothesis's consequences through the
code to localize the suspect set & high \\
\textbf{Induction} & Tests survivors with cheap read-only experiments
(prints, intermediate data) & moderate \\
\end{longtable}
}

\texttt{implement} then writes the surviving hypothesis, with an
adversarial challenger critiquing the diff against the spec.
\texttt{attest} runs the test suite, takes the grader's pass/fail
verdict, and emits a re-entry route (\texttt{inquire},
\texttt{implement}, or \texttt{none}) from a fixed
verdict\(\rightarrow\)route table. The driver parses the verdict and the
route; both are mechanical, and no model decides termination.

\subsection{Hypothesis graph output}\label{recon-output}

\texttt{inquire} emits the hypothesis graph: the structured-analysis
document that precedes the patch. Kill conditions are mechanical
predicates over the evidence trajectory, so a node dies when its
predicate fires and not before. The graph persists across iterations;
re-entry adds nodes rather than overwriting. The frontier closes only
when every open hypothesis is killed (a test refutes it) or witnessed (a
test confirms it).

A committed node is a conclusion, and an inquiry that reaches one rarely
runs straight. Following the \texttt{inquire} skill on a real bug, a
single hypothesis flips across all three modes and a kill before it
settles:

\begin{quote}
abduction \(\rightarrow\) deduction \(\rightarrow\) kill \(\rightarrow\)
abduction \(\rightarrow\) deduction \(\rightarrow\) induction
\(\rightarrow\) deduction \(\rightarrow\) induction \(\Rightarrow\)
induction
\end{quote}

\emph{An in-flight inquiry trace, illustrative: Sonnet 4.6 following the
\texttt{inquire} skill on the python-dotenv \texttt{find\_dotenv} v1.0.1
regression (a real, reproducible bug, every command run). The active
hypothesis cycles through all three modes and a kill before the inquiry
settles; a committed graph records only the terminal node (induction)
and discards this sequence. Full trace:
\texttt{[redacted for double-blind review]}.}

\subsection{Deterministic gating}\label{gating}

The control loop is standard: the driver routes on \texttt{attest}'s
verdict and re-entry route under a bounded attempt budget, and a failure
re-enters \texttt{inquire} with the updated graph rather than retrying
the patch, so the hypothesis graph doubles as the loop's checkpoint and
no dead branch is re-proposed. That leaves the inner layer still owed a
demonstration. The oracle the opening made the crux is what the gate
leans on, and \S\ref{right-regime} isolates it on a single bug, where it
does most of the work the bench number seems to credit to diagnosis.

\subsection{Artifact availability}\label{artifact}

All code and data are openly available, each developed in a public
repository and archived under a DOI for permanence. \texttt{comparator},
the tool that mechanizes the gate, is developed at
\href{https://anonymous.4open.science/r/comparator-DA10}{anonymous.4open.science/r/comparator-DA10}
and archived at
\texttt{[redacted for double-blind review]}
(v1.0.0, AGPL-3.0). The mechanism experiment of \S\ref{right-regime} is
developed at
\href{https://anonymous.4open.science/r/hymech-117E}{anonymous.4open.science/r/hymech-117E}
and archived at
\texttt{[redacted for double-blind review]}
(v1.1.0, CC BY-SA 4.0); the archive carries the preregistration, the
rebuild-confirmed dataset, the regrade script, and the climb traces.
Every reported result reproduces from the committed inputs of the
archived version.

\section{The mechanism by case study}\label{right-regime}

What would prove that we had encoded reasoning into the harness, rather
than just prompted a model into a better answer? The signal has to be
one prompting alone can't produce:

\begin{itemize}
\tightlist
\item
  Without the reasoning operation, the bare model can't reach the fix.
\item
  Handed the operation, it reaches the fix.
\item
  The fix generalizes to cases it was never shown.
\end{itemize}

Here we demonstrate a surprising result on a single bug. Even Fable, the
strongest released model, could not resolve it on its own; the
methodeutic harness carried Sonnet 4.6 to the merged human fix's
behavior on the committed probes.

\subsection{Coding benchmarks measure translation
ability}\label{bench-translation}

We first went looking for that signal on an established coding
benchmark, and learned such benchmarks aren't built to show it. A
SWE-bench-shaped coding benchmark hands the solver three things: the
issue text, the repository at the buggy commit, and a hidden test the
fix must turn from red to green. The issue text often hands over the
specification, removing the need for a discovery step: finding the root
cause and the general fix.

SWE-bench Pro (\href{https://arxiv.org/abs/2509.16941}{Deng et
al.~2025}) is the dominant such benchmark, the one OpenAI now recommends
over Verified. For a competent model, the well-specified tasks are
one-shot from the prompt, so there is no cause to discover. The
underspecified ones are immune to discovery because the author's intent
is hidden from the issue and repository. Our determinacy audit of all
728 public tasks measures this
(\href{https://anonymous.4open.science/r/weaudit-901E}{audit}).

\begin{figure}
\centering
\pandocbounded{\includegraphics[keepaspectratio,alt={SWE-bench Pro's 728 public tasks by determinacy. The green majority is one-shot, with nothing to discover; the 109-task underdetermined floor (15.0\%) grades the author's unstated intent, undiscoverable from the materials.}]{swebench-pro-determinacy.pdf}}
\caption{SWE-bench Pro's 728 public tasks by determinacy. The green
majority is one-shot, with nothing to discover; the 109-task
underdetermined floor (15.0\%) grades the author's unstated intent,
undiscoverable from the materials.}
\end{figure}

With the tests as oracle, the harness resolves 95.3\% of SWE-bench Pro's
public split under the official grader
(\href{https://anonymous.4open.science/r/wepro-3D87}{bench run}): what
iterating against the answer key buys, an artifact of that regime rather
than a leaderboard result, and evidence that the
specification-to-implementation translation step is largely solved. So
the interesting work is now discovery: a bug with a deep, discoverable
cause. Implementation is no longer where the difficulty lives.

Moreover, Pro's public set carried a risk we were not willing to take:
contamination. So we went looking for a post-cutoff bug where discovery
is the whole difficulty. Verus \#2219 is one: a March 2026 issue, opened
and fixed after the solve models' training cutoffs, so the case is
contamination-free. The maintainer's narrow fix (PR \#2230) passes its
shipped test, and a general fix eventually landed too (PR \#2501, merged
2026-06-05). The experiment is openly reproducible on
\href{https://anonymous.4open.science/r/hymech-117E}{GitHub}, with
further analysis within.

\subsection{Verus \#2219}\label{verus-fit}

The bug is
\href{https://github.com/verus-lang/verus/issues/2219}{verus-lang/verus\#2219}:
\emph{ghost uses of never type invalidates borrowchecking}.

Verus is a deductive verifier for Rust: you annotate a program with
specifications and ghost proof code, Verus discharges the obligations
through the Z3 SMT solver, and a passing run certifies the code meets
its spec. The cardinal property of any verifier is \emph{soundness},
that it never certify a program which violates its spec, because every
proof built on top inherits that guarantee. An unsoundness, where the
verifier blesses a program it owes a rejection, is the worst defect it
can carry and the hardest to notice, since the symptom is silence.
Nothing fails.

The correct fix must range over a whole class of program inputs. Two of
them look identical at the \texttt{!} token and owe opposite verdicts:

\begin{Shaded}
\begin{Highlighting}[]
\NormalTok{proof }\KeywordTok{fn}\NormalTok{ unsound(tracked s}\OperatorTok{:}\NormalTok{ S) }\OperatorTok{\{}
\NormalTok{    ghost\_terminal()}\OperatorTok{;}  \CommentTok{// "I STOP here!" ...a lie, proof{-}only, deleted before the program runs,}
\NormalTok{    eat(s)}\OperatorTok{;}            \CommentTok{// so... oops! these two lines actually DO run. Verus thinks they don\textquotesingle{}t,}
\NormalTok{    eat(s)}\OperatorTok{;}            \CommentTok{// but \textasciigrave{}s\textasciigrave{} was already spent. OH NO: bad code waved through. THE }\AlertTok{BUG}\CommentTok{.}
\OperatorTok{\}}

\KeywordTok{fn}\NormalTok{ sound(tracked q}\OperatorTok{:}\NormalTok{ Q) }\OperatorTok{\{}
\NormalTok{    real\_terminal()}\OperatorTok{;}             \CommentTok{// this one really does STOP here (CRASH!)}
\NormalTok{    proof }\OperatorTok{\{}\NormalTok{ take(q)}\OperatorTok{;}\NormalTok{ take(q)}\OperatorTok{;} \OperatorTok{\}}  \CommentTok{// so these genuinely never run, safe to skip. All good!}
\OperatorTok{\}}
\end{Highlighting}
\end{Shaded}

Identical at the token, opposite at the root: the contrast the
experiment needs. The burden of telling them apart falls upstream, which
makes the check stateful and forces it to be exact. Other reasons it is
a good bug to study:

\begin{itemize}
\tightlist
\item
  The merged general fix took deep Verus knowledge and landed three
  months after the issue opened. Not trivial.
\item
  The symptom sits far from the cause: from the wrongly-accepted
  program, the fault sits four steps up the chain (\S\ref{verus-bug}),
  so localization has to climb to find it.
\item
  The project's own suite passes for \emph{both} the narrow and the
  general fix, so it cannot see the distinction the fix has to make.
\end{itemize}

\subsection{Diagnosis}\label{verus-bug}

A \emph{ghost} expression of type \texttt{!}, erased before compilation
and so not actually diverging, still triggers rustc's never-type edge
prune: Verus drops a live control-flow edge, skips the reachable code
behind it, and verifies a program it owes a rejection. That is the
unsoundness.

\begin{figure}
\centering
\pandocbounded{\includegraphics[keepaspectratio,alt={The \#2219 unsoundness as a causal chain. The axis at the bottom is the crux: a ghost-erased uninhabited return (keep the edge) and genuine runtime divergence (prune) are identical at the token, opposite in handling.}]{verus-2219-unsoundness.pdf}}
\caption{The \#2219 unsoundness as a causal chain. The axis at the
bottom is the crux: a ghost-erased uninhabited return (keep the edge)
and genuine runtime divergence (prune) are identical at the token,
opposite in handling.}
\end{figure}

The narrow fix keys on the surface token: when the diverging expression
is a literal \texttt{!}, keep the edge. The maintainer's first patch
(\#2230) does exactly this, and it is the ceiling every self-graded
method reaches. The general fix keeps the edge for the whole class of
uninhabited types the \texttt{!} token never names, which a
surface-token fix cannot reach. That distinction is the XOR the
experiment lives on.

That diagnosis was not handed down; an inquiry climbed to it, and the
climb is retained:

\begin{quote}
\textbf{H0 \textperiodcentered{} observation} \emph{(induction)}.
\texttt{reported.rs} rejected pre-patch with Rust E0382, ``use of moved
value''. \emph{Edge:} is E0382 incidental? Vary the tracked type's
copyability and the never-expression shape.

\textbf{H1 \textperiodcentered{} the erased \texttt{!} is load-bearing}
\emph{(abduction)}. Trial \texttt{copy\_int\_never.rs}: erased
\texttt{!} before a double-used Copy int. \emph{Observed:} ``3 verified,
0 errors'', then a rustc ICE, ``bad arg \#0 (\texttt{!} \(\leftarrow\)
\texttt{()})''. \emph{Verdict:} the erased \texttt{!} corrupts MIR even
where E0382 cannot fire. \emph{Trajectory:} convergent on the erased
ghost never-value.

\textbf{H2 \textperiodcentered{} the source expression is incidental}
\emph{(abduction)}. Trial \texttt{never\_let\_binding.rs}:
\texttt{let\ x:\ !\ =\ arbitrary();\ test(x)}. \emph{Observed:}
rejected, ``never-to-any coercion not allowed in spec mode''.
\emph{Verdict:} the explicit-coercion path is already covered; the
uncovered path is the direct erased call in argument position.
\emph{Kill:} ``all never coercions are uncovered.''

\textbf{H3 \textperiodcentered{} root cause} \emph{(deduction)}. Spec
calls erase to \texttt{erased\_ghost\_value} at the original Rust type;
when that type is \texttt{!}, the injected verifier-only value is
never-typed and breaks the MIR path.

\textbf{H4 \textperiodcentered{} fix} \emph{(deduction)}. Reject erased
ghost calls whose adjusted Rust type is \texttt{!}, before erasure
reaches MIR construction.
\end{quote}

\emph{One inquiry into \#2219 (Sonnet 4.6), condensed from the full
record. It reached the surface-token narrow fix (the \#2230 ceiling) and
stalled at a candidate it could not rebuild to verify. One markdown file
per inquiry; the runs left many:
\href{https://anonymous.4open.science/r/hymech-117E}{pilot
11}.}

\subsection{Experiment design}\label{verus-design}

To demonstrate the mechanism, we hold one reproducible control and
ablate the factors for attribution. Estimating how often the mechanism
is responsible for the outcome is a separate, statistical question, out
of scope here (\S\ref{null-regime} bounds where it engages).

Here, we demonstrate two ablations of interest:

\begin{itemize}
\tightlist
\item
  \textbf{Verification mechanism.} How each kill gets its verdict: by
  the model's own attestation (\emph{self-attested}) or against a
  reference the model cannot author (\emph{externally verified}).
  External means external to the model; this axis is defined in
  \S\ref{method}.
\item
  \textbf{Methodeutic inquiry.} Whether the run runs the typed
  three-mode loop, and at which tier abduction is applied: written into
  the prompt, or lifted out as a tool call (\texttt{comparator}). The
  baseline leaves it to a bare prompt with neither.
\end{itemize}

Six methods self-attest, one is externally verified, ordered from least
to most intervention, with minimal and neutral sharing a row (identical
settings, both run). \emph{Enumeration} here is breadth of search rather
than one guess, over suspect code sites (the \emph{site-enumeration}
arm) or a wide space of discriminating test programs (\hvfull{} present,
\hvempty{} absent):

{\def\LTcaptype{none} % do not increment counter
\begin{longtable}[]{@{}lccc@{}}
\toprule\noalign{}
Experiment arm & Enumeration & Kill conditions & Externally verified \\
\midrule\noalign{}
\endhead
\bottomrule\noalign{}
\endlastfoot
\emph{minimal/neutral} prompt & \hvempty{} & \hvempty{} & \hvempty{} \\
\emph{site-enumeration} prompt & \hvfull{} & \hvempty{} & \hvempty{} \\
\emph{abduction} prompt & \hvempty{} & \hvfull{} & \hvempty{} \\
\emph{hypothesis graph} prompt & \hvempty{} & \hvfull{} & \hvempty{} \\
\emph{self-verifier} harness & \hvfull{} & \hvfull{} & \hvempty{} \\
\emph{comparator}-enabled harness & \hvfull{} & \hvfull{} & \hvfull{} \\
\end{longtable}
}

The \emph{abduction} and \emph{hypothesis graph} prompt arms carry
identical factor settings above; they differ only in whether the run
persists the structured graph across iterations, which the results
(\S\ref{verus}) show makes no difference here. That the graph itself is
not the active ingredient is the point: what moves the outcome is the
verdict source, not the structure several arms share.

Each arm writes a hypothesis graph by the same loop and runs three
times. (Oracle here is metrological: an instrument checked against a
reference, distinct from the probability-calibration of confidence
scores.) We drew candidates from the localization-hard band the lead
case sits in. Before any arm runs, the protocol preregisters one
sentence per loop: testing X, predict Y, refuted by Z.

For ecological accuracy, each model resides in its own native agent
rather than a uniform rig we impose: Claude models in Claude Code,
GPT-5.5 in codex, Composer 2.5 in Cursor. Our harness is a thin wrapper
that drives each vendor's CLI non-interactively with the same stage
prompt:

\begin{Shaded}
\begin{Highlighting}[]
\ExtensionTok{claude} \AttributeTok{{-}p}\NormalTok{ …       }\CommentTok{\# Claude Code}
\ExtensionTok{codex}\NormalTok{ exec …      }\CommentTok{\# codex}
\ExtensionTok{cursor{-}agent}\NormalTok{ …    }\CommentTok{\# Cursor}
\end{Highlighting}
\end{Shaded}

The meta-loop owns only the stage contracts, the hypothesis graph, and
the gate; every model-facing call goes through the vendor's own agent.
Full configs, prompts, and per-arm logs live in
\href{https://anonymous.4open.science/r/hymech-117E}{the mechanism
repository}.

But external verification is necessary. It is not sufficient. A
reference only helps as far as its cases cover the XOR the fix must draw
(\S\ref{verus-bug}), so what matters next is where those correct
verdicts come from and which side of the distinction each one reaches.

\subsection{Making a tiny bench}\label{verus-bench}

To attribute the lift to the mechanism, every arm has to be scored on
the same mechanical ruler; otherwise ``more general'' stays an
impression no one can check. The project's own suite is no such ruler:
it passes for both the narrow and the general fix (\S\ref{verus-fit}).
So we build one by hand, a tiny answer key. The case is one bug, so the
bench is one bug. It holds three compiler states, the buggy commit
(base) and the maintainer's narrow and general fixes, plus a handful of
small test programs, the \emph{probes}. A correct compiler owes each
probe one \emph{correct verdict}: \texttt{REJECT} for an unsound
program, \texttt{ACCEPT} for a sound one. Grading a fix is then
mechanical: does it return each probe's correct verdict?

{\def\LTcaptype{none} % do not increment counter
\begin{longtable}[]{@{}lcl@{}}
\toprule\noalign{}
Probe program & Golden assessment & Who supplied it \\
\midrule\noalign{}
\endhead
\bottomrule\noalign{}
\endlastfoot
bug probes (uninhabited) & \texttt{REJECT} & base, by construction
(free) \\
divergence, in-bar & \texttt{ACCEPT} & the merged human fix (\#2501) \\
\end{longtable}
}

The split is the whole asymmetry. Bug probes are uninhabited \emph{by
construction}, so their correct \texttt{REJECT} is free and base is
wrong on exactly them. The gate (\S\ref{gate-general}) enumerates such
cases and passes a fix only when it flips them all to \texttt{REJECT}
with no sound case newly rejected, no human in the loop. Divergence
probes are the hard side, sound and to be preserved, but base is no
reference there and the gate's grammar holds no divergence-preserve
shape, so the gate is blind to them. Their \texttt{ACCEPT} comes from a
human, the merged fix \#2501 (with one stretch case even \#2501
declines), and stays held out of the gate. Passing it therefore tests
whether a fix represents the rule rather than fits the gate. The easy
side's verdict is free; the hard side's costs a human
(\S\ref{enum-calib}).

A recall probe rules out memorization: asked in isolation how \#2219 was
fixed, the model does not recover the patch, so its solution is
reconstruction. Every verdict is from a forced-fresh, fingerprinted
build; the frozen dataset and regrade script are committed.

\subsection{Results}\label{verus}

Which arm climbs past the narrow fix? The six self-attested methods
plateaued on all eighteen draws, three each; only the
externally-verified arm broke past.

{\def\LTcaptype{none} % do not increment counter
\begin{longtable}[]{@{}
  >{\raggedright\arraybackslash}p{(\linewidth - 6\tabcolsep) * \real{0.2500}}
  >{\centering\arraybackslash}p{(\linewidth - 6\tabcolsep) * \real{0.2500}}
  >{\centering\arraybackslash}p{(\linewidth - 6\tabcolsep) * \real{0.2500}}
  >{\raggedright\arraybackslash}p{(\linewidth - 6\tabcolsep) * \real{0.2500}}@{}}
\toprule\noalign{}
\begin{minipage}[b]{\linewidth}\raggedright
Arm
\end{minipage} & \begin{minipage}[b]{\linewidth}\centering
Gate pass
\end{minipage} & \begin{minipage}[b]{\linewidth}\centering
Cases flipped
\end{minipage} & \begin{minipage}[b]{\linewidth}\raggedright
Outcome
\end{minipage} \\
\midrule\noalign{}
\endhead
\bottomrule\noalign{}
\endlastfoot
minimal / neutral & \xmark{} & 114 & narrow plateau (\#2230 slice) \\
site-enumeration & \xmark{} & 114 & narrow plateau \\
abduction & \xmark{} & 114 & narrow plateau \\
hypothesis graph & \xmark{} & 114 & narrow plateau \\
self-verifier & \xmark{} & 114 & narrow plateau \\
\texttt{comparator} & \cmark{} & 269 & general on the bug-set,
wide-but-broken on divergence \\
\end{longtable}
}

\emph{Cases flipped is the modal of three draws; gate pass is the gate's
verdict on its own cases.}

The six self-attested methods all stopped at the narrow plateau, none
reaching \texttt{pass=true}. Only the \texttt{comparator} arm broke off
it, with zero valid-preserve rejections on the gate's own cases. Its fix
is more general than the plateau but not yet the merged human fix: it
stays wide-but-broken on the in-bar divergence case the gate never
enumerated (\S\ref{gate-general}). The reach is real: it flips the whole
bug-set and rejects two \emph{out-of-grammar} held-outs the gate never
showed it, applying its general predicate to cases it never saw.

That isolates the mechanism on one model. The headline lift is a slice
across two: the same bug handed to the strongest released model without
the comparator, and to a weaker one with the comparator on a gate whose
coverage includes the divergence case (\S\ref{enum-calib}); a gate
silent on divergence leaves both arms wide-but-broken, as above.

{\def\LTcaptype{none} % do not increment counter
\begin{longtable}[]{@{}lcc@{}}
\toprule\noalign{}
& No comparator & + comparator (gate covering divergence) \\
\midrule\noalign{}
\endhead
\bottomrule\noalign{}
\endlastfoot
\textbf{Fable 5} (strongest released) & wide-but-broken & \#2501
behavior \\
\textbf{Sonnet 4.6} & narrow (\#2230) & \#2501 behavior \\
\end{longtable}
}

\emph{Two contamination-clean models, one harness, the gate covering
divergence (\S\ref{enum-calib}). Without the comparator neither matches
the merged fix, and they miss in opposite directions: Fable too wide,
Sonnet too narrow. With it both reach \#2501's behavior on the committed
probes (including \#2501's own over-conservatism on the stretch case),
which the stronger model cannot reach on its own. The comparator returns
only pass/fail on candidate behavior, never the patch or the predicate
(\S\ref{enum-calib}), so the match is reconstruction, not lookup. One
draw per cell.}

\subsection{\texorpdfstring{\texttt{comparator}}{comparator}}\label{gate-general}

\texttt{comparator} is an open-source command-line tool anyone can run
(\href{https://anonymous.4open.science/r/comparator-DA10}{anonymous.4open.science/r/comparator-DA10}):
the diff of \S\ref{application}, the surprise-manufacture, built as a
standalone instrument on the program-analysis lineage that put
counterexample search to work (\S\ref{lineage}). The name is legacy;
what it does is provoke, not abduce. On this bug it carried a model past
the plateau prompting alone could not clear. Our read of why is that it
externalizes that surprise from the inquiry loop into three
domain-general operations with no answer built in:

\begin{itemize}
\tightlist
\item
  It \emph{enumerates} a space of cases closed under the property's
  type-formers, wider than any one hypothesis.
\item
  It \emph{calibrates} each case against a known-good baseline, the
  comparator the model cannot author, so the ground truth is external to
  it.
\item
  It \emph{gates} on a single pass/fail signal the model iterates
  against.
\end{itemize}

The prompt demands generality and supplies the signal, so the rule is
the model's own reconstruction. Handed only the instruction to range
over type-formers and judge, Fable \emph{rebuilt \texttt{comparator}
itself} from the goal alone (\S\ref{enum-calib}): its own 7,026-case
gate over the uninhabited type-formers. A general construction permits
that; a bespoke one, with the answer baked in, would not.

The gate's coverage is the frontier, lever and limit at once. Where it
covers, the uninhabited side, the model greps the codebase and widens
its predicate to clear the cases the gate feeds, supplying the discovery
while the gate gives direction. Where its grammar is silent, the
divergence side, it over-generalizes: the fix lands general on the bug
set and over-conservative on divergence, \emph{wide but broken}, the
\S\ref{verus-bug} XOR collapsed to an OR where the gate stayed silent.
The lift it buys is a coarse mode gate. It is no recovery of the
verifier's decision procedure: a fix keeping the edge for every
ghost-mode call, with no inhabitedness query at all, grades identically
on every probe. Supply the missing divergence verdicts
(\S\ref{enum-calib}) and the arm reaches the human fix.

The construction generalizes past this bug, and past verification: the
same three operations reconcile any two accept-sets, a refactor against
its original or a port against its source. They recover the disagreement
from a compact sketch, so neither set need enter a context window or
cross a wire. The full command surface and the cited lineage are in the
\href{https://anonymous.4open.science/r/comparator-DA10}{repository}.

\subsection{The surprise is not self-generable}\label{enum-calib}

A model can bootstrap the \emph{enumeration}, the combinatorial breadth,
because that is mechanical: handed a vague prompt and no gate, Fable
rebuilt one for itself (\S\ref{gate-general}). It cannot bootstrap the
\emph{correct verdicts}, the human-accepted answers. Its self-built gate
labels each case with the very predicate under test, so the one case
that needs a truth from outside its own belief, genuine divergence
versus ghost-erasure, gets self-mislabeled as handled. One run shows
both halves: a wide net the model built, and a blind spot exactly where
its own labels could not reach. The claim is scoped to that self-grading
circularity. It is not a proof that no model could ever produce a
correct verdict some other way.

So before crediting the verdict source, the cheaper explanations have to
fall, and each has a control already in the ablation.

\emph{Did the comparator just search harder?} The self-verifier searched
exactly as hard: handed the same prompt and no tool, it built its own
case generator and brute-forced some 7,026 cases a round, the breadth
the comparator enumerates, and it still plateaued: confirmation bias
(\S\ref{grounding}) with a receipt, grading against its own belief, the
one move the external comparator forbids. Enumeration sat on both sides
of the comparison, so it is not the lever. Nor is reasoning depth: the
abduction arm ran textbook bi-abduction to the semantic predicate the
general fix turns on, named it in prose, self-validated it against its
own programs, and still plateaued at the same narrow verdict as the
one-line minimal arm. An arm can write down the correct predicate and
certify it against its own belief; the patch ships narrow all the same.
The two-model rerun sharpens the point: without the comparator the models
miss in opposite directions, Fable too wide and Sonnet too narrow, and
the comparator pulls both to the same \#2501. More search does not correct
two opposite errors onto one target; an external answer key does.

\emph{Was the model handed the fix?} It was handed correct verdicts on
cases. It never saw the patch, the predicate, or the Rust API. The
comparator returns only pass/fail on candidate behavior, so the model
still localizes the fault four steps up the chain (\S\ref{verus-bug})
and writes the fix itself. This is why the signal is verdicts and not
the issue text: where a benchmark leans on the issue to say what
\emph{fixed} means, and hands over the spec doing it, the gate says it
in correct verdicts on cases, the more specific signal, grading whether
the fix draws the exact behavioral boundary rather than whether it
parsed a prose description.

\emph{A better scaffold, then?} The self-verifier is the scaffold
control: the whole loop, enumeration and kill conditions included, with
the external verdict the only piece removed. It plateaus beside the bare
prompt. A scaffold without the answer key buys nothing here.

\emph{Just the stronger model?} The attribution rests on the
within-model contrast, where the model is held fixed and only the
verdict source moves. The cross-model grid (\S\ref{verus}) only
illustrates the consequence; the within-model ablation is what proves
the cause. It shows the lift is wide enough that a weaker model with the
verdicts passes a stronger one without them.

The whole ablation reduces to one causal diagram: across its arms,
everything stays fixed but the verdict source, the one input the model
cannot author.

\begin{figure}
\centering
\pandocbounded{\includegraphics[keepaspectratio,alt={Reading the diagram: each box is an action, read left to right. Four factors stay identical across the ablation's arms (model, loop, bug, graph), so none of them can be the cause; only the verdict source varies. Gray boxes the model can build for itself; the blue box is the one input it cannot author, and the blue path is the causal route. The gray arrow into the outcome is the implementation caveat, where codex walls.}]{verus-2219-lift-mechanism.pdf}}
\caption{Reading the diagram: each box is an action, read left to right.
Four factors stay identical across the ablation's arms (model, loop,
bug, graph), so none of them can be the cause; only the verdict source
varies. Gray boxes the model can build for itself; the blue box is the
one input it cannot author, and the blue path is the causal route. The
gray arrow into the outcome is the implementation caveat, where codex
walls.}
\end{figure}

\emph{Not one model's artifact.} The \S\ref{verus} lift is one model
pair; extend the gate's coverage to the divergence case, verifying
against the correct verdicts (\#2501), and the same fix lands across
four models in four native CLIs:

{\def\LTcaptype{none} % do not increment counter
\begin{longtable}[]{@{}lcc@{}}
\toprule\noalign{}
Verdict source & Narrow fix (\#2230) & General fix (\#2501) \\
\midrule\noalign{}
\endhead
\bottomrule\noalign{}
\endlastfoot
self-attested, six methods & \hvfull{} & \hvempty{} \\
externally verified --- \emph{Fable} & \hvfull{} & \hvfull{} \\
externally verified --- \emph{Sonnet 4.6} & \hvfull{} & \hvfull{} \\
externally verified --- \emph{Composer 2.5} & \hvfull{} & \hvfull{} \\
externally verified --- \emph{codex} & \hvfull{} & \hvempty{} \\
\end{longtable}
}

codex clears the bug arm but cannot implement the divergence case, an
implementation limit the verdicts do not remove.

The convergence is coverage-bound: it shows the fix is reproducible
across workflows once the gate covers divergence. It does not show three
models independently rediscovered the predicate. The shared decline on
the stretch case is most likely the gate funnelling every successful arm
into the one behavior it rewards, which the human fix happens to share.

Retrospectively the correct verdicts are free; prospectively, on a fresh
bug with no merged fix, the easy side has one and the hard side does
not. The division of labor is the deployment design: the model
enumerates and fixes, the harness draws its verdicts from approved
history, and the residual hard side gets named. Nothing is hidden.

\subsection{The surprise becomes a tool call}\label{abduction-tool}

Step back from the arms. What the traces catch is the models computing
the diff by hand: the self-verifier recomputed the full symmetric
difference every pass (\S\ref{enum-calib}), the tell that the XOR is the
surprise abduction runs on, its input and not the leap. But a model
grading its own cases can only check them against its own belief, so
nothing surprises it, and it plateaued. The model could make the leap
all along; what it lacked was a surprise it could not grade away, one
from outside its own belief.

That is the signal we went looking for. What lifts off the model is not
abduction but the surprise that triggers it: the symmetric difference
between what the model believes and what is true is a counterexample,
and in coding diagnosis it runs as a tool call, the one we release as
\texttt{comparator}: enumerate, calibrate, gate. The leap that answers the
surprise stays in the model, uncoded, because everything we managed to
mechanize turned out to be the test or the record, never the guess. That
division, the harness manufacturing the surprise and the model making
the leap, is the division of intellect the introduction set out to buy,
and the run leaves behind the hypothesis graph the introduction
promised, written by the experiment that needed it.

\section{Discussion}\label{discussion}

\subsection{Diagnostic frontier}\label{null-regime}

Implementation is the solved part. A competent model turns a
well-specified issue into a passing patch, which is most of what a
SWE-bench task rewards (\S\ref{bench-translation}) and most of what
coding agents already do. What is left is \emph{discovery}: the root
cause when the symptom sits far from it, the general fix when the suite
cannot tell it from a narrow one. That diagnostic step is the frontier,
and it is where the mechanism engages.

The band is narrow precisely because implementation is solved, so most
bugs never reach the diagnostic step. On pilots from the merged-PR pool
and beyond, a minimal baseline solved them unaided, from trivial fixes
(qrtool \#695, bat \#3710) to the cases that most invited the machinery
(slang-server \#310, fjall \#287). Two findings attribute the nulls. A
\textbf{selection artifact}: the deployment pipeline's own triage
fast-paths easy bugs around the graph (``a 1-line fix with a confirmed
reproducer doesn't need a hypothesis graph''), so the merged pool is
exactly the implementation-solved subset. And \textbf{baseline reach}: a
capable model in a minimal loop fixes most reproducible bugs unaided.
The band is bounded below by that reach and above by triage, which is
why a population rate over an unselected pool measures the band's width
and misses the mechanism; the Verus case is chosen to sit inside it.
These nulls publish \emph{attributed}, each arriving with the mechanism
that produced it, so an unexplained null says stop where an attributed
one says where to point the next instrument.

Automation is a different axis. A faster loop, a parallel fan-out, a
cheaper scaffold moves the same work along quicker, and none of it
changes whether the agent can find the cause. The lift here is on the
diagnostic axis, the distinction the scaffold-synergy results blur
(\S\ref{rw-scaffolds}): a weaker model with better plumbing finishing
sooner is a throughput gain, where a weaker model reaching a fix a
stronger one could not is a capability gain.

\subsection{Verifiable agents}\label{verifiable-agents}

The reframe is from treating an agent as a source of verdicts to a
source of claims bound to evidence.

The failure mode is output that is confidently wrong and not cheaply
verifiable, and it scales with model capability rather than against it,
the default product of fluent generation. The Verus self-built gate is
the instance (\S\ref{enum-calib}): green by the agent's own check and
unsound at once, exposed only by the held-out probes replayed against
the build. The same shape recurred throughout, in the wide-but-broken
fix, the 4/4 false-green self-audit, and the recall-inflated resolve
rate.

The substitution is accountability in place of trust, checkable line by
line. The self-graded agent produced a verdict, its own gate green; the
externally-oracled inquiry produced a record of typed nodes, each a
hypothesis, an exact command, an observed outcome, and the edge the
result generates, every consequential one replaying on a clean build.
Such a record is audited step by step without extending trust to its
author, and the better fix is identifiable independent of the system
that produced it. That property, rather than the outcome of any single
comparison, is what survives.

\textbf{Truth is buildable.} The node semantics of \S\ref{epistemics} is
what gives the trace its force here. If a true claim is a build
presently passing and its warrant lives in the edges, then a node
without a replayable trial is not a node at all, as an uncheckable
number is not a measurement. A resolve rate, however honest, is a
verdict over 728 builds an auditor cannot inspect; a hypothesis graph is
the inspectable chain.

\textbf{The cost asymmetry of fabrication.} A fabricated reasoning trace
is expensive to sustain, because every fabricated node must survive a
replay the author does not control. A confident narrative is cheap to
produce, because nothing in it is bound to a procedure. A self-graded
method's over-narrow fix is the cheap kind: it reads as complete, and
closer reading does not reveal the flaw, while the receipt reveals it in
one command. Verification is therefore part of the method rather than
overhead on it. The same asymmetry governs failure: a trusted verdict
that fails leaves nothing behind, whereas a recorded inquiry that fails
leaves a trail naming the failed node. The Verus climb recorded its own
corrections in trail, the narrow \texttt{is\_never} fix refuted by the
very cases the gate fed it, each a kill that generated the next edge and
itself a replayable trial.

\textbf{Hidden effort is not reasoning you can check.} The frontier labs
ship ``reasoning'' as an opaque dial (\texttt{high}, \texttt{xhigh},
\texttt{ultrathink}): more private tokens burned before the answer, the
model grading against its own belief, harder. The Verus result is where
that runs out (\S\ref{enum-calib}). Effort scales the inducible half,
the enumeration, and cannot manufacture the half that is not, a
ground-truth oracle from outside the model's belief, so the
self-verifier plateaus where a single externally-graded kill reaches the
general fix. Whatever the dial buys ships nothing a stranger can rerun;
an effort setting is the purest \emph{take my word for it}, a larger
number where a kill condition should be. The position here is the
opposite: put the reasoning in the harness, typed and replayable, its
level a trail the auditor checks instead of a knob the auditor trusts.

\textbf{Accountability, in deployment.} The intro's promise was a claim
that clears only when an independent party reruns its recorded trial
(\S\ref{introduction}); the receipts deliver it. To a maintainer the
agent and its operator are strangers, and with the trial attached that
stops mattering: a red-on-master/green-with-fix test and a clean suite
read identically whoever submitted them. This is the review bottleneck
answered by making the check cheaper than reconstructing the reasoning.

\textbf{Merit attaches to the work, not the doer.} An agent can produce
many artifacts of any quality, so judging them by their author is
backwards; judging the author by replaying the artifacts is the only
direction that scales. Merit is the warrant a piece of work carries in
itself, checkable without reference to who or what produced it: the
hypothesis graph is a unit of work that ships with its own evidence, the
receipts-first PR the corresponding practice, the maintainer who merges
on the ledger alone an early instance.

\section{Related work}\label{related-work}

\subsection{Construct validity and contamination}\label{rw-swebench}

The SWE-bench family defines the Verified / Pro lineage, official
harness, and contamination-resistant tier design. \textbf{SWE-Bench+}
(Aleithan et al.~2024) manually audited the original bench: 32.67\%
solution leakage, 31\% weak tests. \textbf{OpenAI's February 2026 audit}
found a majority of audited Verified tasks have flawed tests and that
frontier models reproduce exact gold patches; it stopped reporting
Verified and recommends Pro. \textbf{Wang, Pradel \& Liu} (ICSE 2026)
show plausible patches pass tests yet diverge from developer intent;
their axis is patches that pass but are wrong, ours
(\S\ref{right-regime}) is tasks whose materials do not determine which
passing behavior is intended. \textbf{ORACLE-SWE}
(\href{https://arxiv.org/abs/2604.07789}{arXiv:2604.07789}) quantifies
the same handover, ablating the oracle and specification signals that
leak through a task and measuring the resulting drop; \textbf{SLUMP}
(\href{https://arxiv.org/abs/2603.17104}{arXiv:2603.17104}) opens on the
identical premise, that benchmarks supply the full specification upfront
while real coding does not, and answers it by building an
underspecified-by-design benchmark. Here we part from both on the
verdict: they treat handover as a defect to fix with a better benchmark,
while the determinacy audit (\S\ref{right-regime}) draws it as a
category boundary, a spec-conformance instrument cannot be tuned into a
measure of diagnostic inquiry because the two are different types.
Evaluation surveys reach the same edge from the methodology side:
scoring final outputs misses the in-run reasoning and failure causes,
and trajectory-level assessment is the called-for remedy (Yehudai et
al.~2025), which is what grading the trace rather than the patch
(\S\ref{future-work}) would supply. \textbf{SWE-rebench} uses
post-cutoff filtering as a parallel contamination strategy;
\textbf{LiveCodeBench} (Jain et al.~2024) is the origin of post-cutoff
(temporal-holdout) evaluation, and the standard objection to it applies
here too, that training cutoffs are porous because RL post-training and
inference-time retrieval can surface later content. The witnessed case
(\S\ref{verus-bench}) is built against exactly that objection: it is
post-cutoff and the solving model's weights predate the fix, so neither
porosity nor retrieval supplies the answer. \textbf{HAL} (Stroebl et
al.~2025) sets the cost-transparency precedent and the official
\textbf{swe-bench/experiments} repo the per-instance publication norm
(\texttt{trajs}/\texttt{logs}/\texttt{patch}/\texttt{report}); the
receipts here extend both with gate traces, hypothesis graphs, and a
re-gradeable cost ledger.

\subsection{Agent scaffolds and SE-agent harnesses}\label{rw-scaffolds}

Surveys and position papers now map the harness's responsibilities
directly: the case for \emph{programming with trust}, verification and
analysis built into the agent rather than retrofitted after
(Roychoudhury et al.~2025), the systematizing of agents for software
engineering around verification, testing, and repair (Liu et al.~2024),
and persistent, structured memory named among the open challenges for
long-running agents (Wang et al.~2025). We answer with the structure
none of them name: a Peirce-typed, kill-conditioned hypothesis-graph
memory. The SWE-bench-targeted harnesses that exist, OpenHands (Wang et
al.~2024), SWE-agent (Yang et al.~2024), and AutoCodeRover (Zhang et
al.~2024/25), are ReAct-pattern loops (Yao et al.~2023) without it.
\textbf{Voyager} (Wang et al.~2023) is the closest loop-shape precedent:
embodied
observe\(\rightarrow\)hypothesize\(\rightarrow\)test\(\rightarrow\)commit,
with a skill library where this work holds falsifiable claims.
\textbf{SWE-Effi}
(\href{https://arxiv.org/abs/2509.09853}{arXiv:2509.09853}) is the
sharpest published counter-position: effectiveness emerges from
scaffold-model synergy rather than residing in the scaffold alone. Here
we agree from the other direction, with the synergy named: the binding
pair is gate \(\times\) oracle, isolated in the controlled experiment of
\S\ref{right-regime} rather than on a benchmark that cannot see it.

Recent results do report a weaker model with a scaffold beating a
stronger one, on the automation axis rather than the capability one.
\textbf{GradleFixer}
(\href{https://arxiv.org/abs/2510.08640}{arXiv:2510.08640}) has
Gemini-2.5-Flash out-repair Gemini-2.5-Pro on Android builds, one
generation apart inside a single family; the lift is
\emph{domain-specific} execution primitives for Gradle, so a model
handed purpose-built tools for one build system doing better is
expected. \textbf{Confucius Code Agent}
(\href{https://arxiv.org/abs/2512.10398}{arXiv:2512.10398}) has Sonnet
4.5 edge Opus 4.5 on SWE-bench Pro, 52.7 to 52.0, with the edge
reversing once Opus runs the same scaffold (54.3), a scaffold effect
rather than a capability crossing. Both externalize actions or notes;
the diagnosis stays in the model.

The lift in \S\ref{right-regime} parts from these on three counts: it is
a frontier extension rather than a throughput gain; the externalized
operation is \emph{domain-general} (the released \texttt{comparator}
carries no answer and no domain model, \S\ref{gate-general}) rather than
a tool shaped to one build system; and a controlled ablation pins it to
the verdict source alone, so it does not reverse when the stronger model
is handed the same harness, the way a whole-scaffold swap does. The
nearest comparator precedent, \textbf{SWT-Bench}
(\href{https://arxiv.org/abs/2406.12952}{arXiv:2406.12952}), grades
generated tests by fail-on-original then pass-after-golden-patch, the
same external-golden shape, but as an evaluation oracle for test
generation. Ours is the runtime instrument an agent iterates against
with the answer key hidden.

Two concurrent developments arrived independently at adjacent points,
each carrying one of the two components composed here.
\textbf{Theorem-of-Thought}
(\href{https://arxiv.org/abs/2506.07106}{Abdaljalil et al.~2025}) types
reasoning into abductive, deductive, and inductive specialist agents per
query: the typed cycle, without a persistent typed memory across cycles.
\textbf{Cognitive Memory Manager}
(\href{https://openreview.net/forum?id=yCsHQnvvWY}{Khalid \& Arora
2026}) extracts a typed-node DAG by observing agent execution and mines
it for patterns to promote to skills: the typed graph, mined
descriptively where ours is generative (it routes the run). That
convergence is independent. Provenance for the framing here is
timestamped on the project blog
(\texttt{[redacted for double-blind review]},
\texttt{[redacted for double-blind review]}).

The trichotomy these siblings reach for is Peirce's (1878, 1903).
\S\ref{grounding} and \S\ref{lineage} wire the vocabulary to its
sources, and the dated posts above timestamp this lineage's use of the
hypothesis-graph primitive.

{\def\LTcaptype{none} % do not increment counter
\begin{longtable}[]{@{}
  >{\raggedright\arraybackslash}p{(\linewidth - 8\tabcolsep) * \real{0.1918}}
  >{\raggedright\arraybackslash}p{(\linewidth - 8\tabcolsep) * \real{0.1096}}
  >{\raggedright\arraybackslash}p{(\linewidth - 8\tabcolsep) * \real{0.2466}}
  >{\raggedright\arraybackslash}p{(\linewidth - 8\tabcolsep) * \real{0.2466}}
  >{\raggedright\arraybackslash}p{(\linewidth - 8\tabcolsep) * \real{0.2055}}@{}}
\toprule\noalign{}
\begin{minipage}[b]{\linewidth}\raggedright
System
\end{minipage} & \begin{minipage}[b]{\linewidth}\raggedright
Domain
\end{minipage} & \begin{minipage}[b]{\linewidth}\raggedright
Reasoning-mode typing
\end{minipage} & \begin{minipage}[b]{\linewidth}\raggedright
Persistent structure \& update
\end{minipage} & \begin{minipage}[b]{\linewidth}\raggedright
Termination gate
\end{minipage} \\
\midrule\noalign{}
\endhead
\bottomrule\noalign{}
\endlastfoot
Voyager (Wang et al.~2023) & Minecraft & None & Skill library;
test-validated graduation & Test-pass on skill \\
IDEA (He et al.~2025) & Interactive rule learning & Peirce-cited,
agent-level & Working rule set & None explicit \\
ADI (Gilda \& Gilda 2026) & Algebraic invariants & Peirce, layered
(L0/L1/L2) & Symbolic knowledge graph & None explicit \\
AriGraph (Anokhin et al.~2024) & TextWorld & None & Knowledge graph
(entities, relations, episodes) & None explicit \\
CausaLab (Yang et al.~2026) & Causal discovery & Causal-typed (SCM) &
Evolving structural causal model in a DSL & None explicit \\
BeliefMem (Liao et al.~2026) & Partial-observability QA & None &
Candidate set; Noisy-OR probabilistic update & Probabilistic
threshold \\
Theorem-of-Thought (Abdaljalil et al.~2025) & General reasoning &
Abduction / deduction / induction, agent-level & Formal reasoning graph
& NLI-guided Bayesian coherence \\
CMM (Khalid \& Arora 2026) & SE (coding agents) & 7 trajectory roles,
extraction-time & Typed DAG; confidence decay & Human approval +
retrieval-validated threshold \\
This work & SE (industrial code) & Peirce, enforced at write time per
stage & Hypothesis graph; mechanical kill predicates on the audit
verdict & Deterministic finite-state \\
\end{longtable}
}

\emph{Table 1. Comparison spine for adjacent typed-reasoning and
graph-memory LLM-agent systems. Cell terseness is by design; prose
nuance in \S\ref{typed-memory}.}

\subsection{Typed reasoning and graph-structured
memory}\label{typed-memory}

\textbf{IDEA} (He et al.~2025, ACL Findings,
\href{https://arxiv.org/abs/2408.10455}{arXiv:2408.10455}) explicitly
cites Peirce and uses the three modes in an interactive rule-learning
benchmark. \textbf{ADI} (Gilda \& Gilda 2026,
\href{https://arxiv.org/abs/2604.15727}{arXiv:2604.15727}) gives an
explicit Peircean tripartite protocol with epistemic layers over a
symbolic knowledge graph; near-simultaneous with this draft and the most
conceptually adjacent prior work. Both target reasoning domains outside
SE.

The hypothesis graph sits at the intersection of three lineages:
cognitive-architecture memory (Soar / ACT-R / EPIC), LLM-agent memory
systems (CoALA / AriGraph / Mem0 / Zep), and typed-belief
representations (CausaLab / BeliefMem / Theorem-of-Thought / CMM). The
hypothesis graph adopts the Soar memory typology directly as its slot
vocabulary, adding only the specific content of the \texttt{smem} slot:
Peirce-typed, kill-conditioned, designed for LLM prose read/write.
Adjacent work: \textbf{Kirk, Wray \& Laird 2023}
(\href{https://ojs.aaai.org/index.php/AAAI-SS/article/download/27690/27463/31741}{AAAI}),
an LLM-port of the Soar lineage; \textbf{CoALA} (Sumers et al.~2023/24,
\href{https://arxiv.org/abs/2309.02427}{arXiv:2309.02427});
\textbf{AriGraph} (Anokhin et al.~2024/25,
\href{https://arxiv.org/abs/2407.04363}{arXiv:2407.04363}), the closest
precedent for graph-structured LLM-agent memory; \textbf{CausaLab} (Yang
et al.~2026, \href{https://arxiv.org/abs/2605.26029}{arXiv:2605.26029});
\textbf{BeliefMem} (Liao et al.~2026,
\href{https://arxiv.org/abs/2605.05583}{arXiv:2605.05583}), strong
adjacent on uncertain alternatives with mechanical update.

\textbf{CMM} (Khalid \& Arora 2026,
\href{https://openreview.net/pdf?id=yCsHQnvvWY}{OpenReview}; a day
before this draft) is the closest comparison: the same persistent typed
DAG of reasoning artifacts, but observe-and-consume (it types a
trajectory post hoc and graduates skills) where ours is
perturb-and-falsify (kills fire live, the graph routes the run). The
directions are opposite and complementary, the \textasciitilde385
committed graphs in \texttt{sweep/repo-hypotheses/} exactly the corpus
its graduation pipeline could consolidate.

Four 2026 systems each carry one component this work combines; what is
new here is the \emph{join}; each piece already exists in one of them.
\textbf{FVDebug}
(\href{https://arxiv.org/abs/2510.15906}{arXiv:2510.15906}) builds an
actual hypothesis graph for debugging, with a frontier and accumulated
evidence, but selects the next node by asking the model, the arbiter
this work removes. \textbf{From Hypotheses to Factors}
(\href{https://arxiv.org/abs/2604.26747}{arXiv:2604.26747}) runs the
same perturb-and-falsify loop, falsifiable hypotheses behind a
deterministic engine over an append-only trace, locked to quantitative
finance where this work claims the general semantic-memory substrate.
\textbf{Portable Agent Memory}
(\href{https://arxiv.org/abs/2605.11032}{arXiv:2605.11032}) is the
nearest provenance memory, a Merkle-DAG that cites the Soar lineage and
makes every node reconstructible by content-addressing, but it certifies
\emph{integrity} (the recorded bytes are untampered) where the replay
invariant here certifies \emph{warrant} (the node still survives its
trial). And the provenance survey \textbf{From Agent Traces to Trust}
(\href{https://arxiv.org/abs/2606.04990}{arXiv:2606.04990}) enumerates
exactly the relations this work mechanizes, Support, Contradict,
Invalidate, and names ``how provenance quality should be evaluated'' as
an open problem; the hypothesis graph is one answer, with replay as the
quality bar and the kill condition as an executable edge rather than a
descriptive label.

What the mechanism experiment (\S\ref{right-regime}) adds to this
comparison is which piece is decisive: not the graph structure, which
several of these share, but where a kill edge gets its ground truth.
FVDebug lets the model arbitrate the next step; the contrast that
separates a narrow fix from a general one is whether the kill is graded
against the model's own belief or against an external oracle, and the
externalized form is released as \texttt{comparator} (enumerate a case
space, calibrate against a known-good reference, gate). A model can
rebuild the enumeration for itself but not that oracle
(\S\ref{enum-calib}), which is the axis these neighbors leave implicit.

A second cluster treats truth and uncertainty as first-class rather than
a downstream score: \textbf{NARS}, \textbf{OpenCog's AtomSpace/PLN},
\textbf{Nanopublications} (Groth et al.~2010), and, closest in time,
\textbf{Traxia}
(\href{https://arxiv.org/abs/2606.08256}{arXiv:2606.08256}), converging
on these primitives two days after \emph{Truth Is Buildable}
(2026-06-04). Where each stops short of a replayable, kill-conditioned
entitlement ledger is adjudicated in \emph{Verifiable Knowledge}, the
paper that owns the epistemology. What is specific \emph{here} is the
data structure: none makes that meaning the \textbf{semantic contract of
a memory node}, truth operationalized by replayable edge structure
rather than a stored label or textual provenance record.

\begin{quote}
\emph{Production LLM memory systems with graph variants (Zep/Graphiti,
Mem0), staged-hypothesis selection in science agents, deterministic
gating in adjacent settings, and reflective memory systems (Reflexion,
DebugMate) are surveyed in the appendix; they are adjacent on particular
axes but do not change the comparison spine.}
\end{quote}

\subsection{Adversarial filtering and
termination}\label{adversarial-termination}

{\def\LTcaptype{none} % do not increment counter
\begin{longtable}[]{@{}
  >{\raggedright\arraybackslash}p{(\linewidth - 8\tabcolsep) * \real{0.3425}}
  >{\raggedright\arraybackslash}p{(\linewidth - 8\tabcolsep) * \real{0.1233}}
  >{\raggedright\arraybackslash}p{(\linewidth - 8\tabcolsep) * \real{0.1644}}
  >{\raggedright\arraybackslash}p{(\linewidth - 8\tabcolsep) * \real{0.2329}}
  >{\raggedright\arraybackslash}p{(\linewidth - 8\tabcolsep) * \real{0.1370}}@{}}
\toprule\noalign{}
\begin{minipage}[b]{\linewidth}\raggedright
System
\end{minipage} & \begin{minipage}[b]{\linewidth}\raggedright
Domain
\end{minipage} & \begin{minipage}[b]{\linewidth}\raggedright
Stage operated at
\end{minipage} & \begin{minipage}[b]{\linewidth}\raggedright
Visibility regime
\end{minipage} & \begin{minipage}[b]{\linewidth}\raggedright
Cross-family
\end{minipage} \\
\midrule\noalign{}
\endhead
\bottomrule\noalign{}
\endlastfoot
Multi-Agent Debate (Liang et al.~2023/24,
\href{https://arxiv.org/abs/2305.19118}{arXiv:2305.19118}) & General
reasoning & Patch / answer stage & Open (cross-visibility) & Single
model family \\
Refute-or-Promote (Agarwal 2026,
\href{https://arxiv.org/abs/2604.19049}{arXiv:2604.19049}) & Defect
discovery & Review stage & Asymmetric context & Yes \\
This work & SE (industrial code) & Pre-patch hypothesis stage & Blind
challenge (no cross-visibility) & Yes (Sonnet + GPT-5.5) \\
\end{longtable}
}

*Table 2. Adversarial multi-model filtering: this work occupies the
pre-patch / blind cell. Termination disciplines (\(\lambda\)\_A's
type-theoretic proofs, SafetyDrift's absorbing states) sit at
composition or trajectory scope where this work's verdict-routed gate
sits per-instance.*

Closest in spirit is \textbf{POPPER}
(\href{https://arxiv.org/abs/2502.09858}{arXiv:2502.09858}), which runs
agentic sequential hypothesis tests under e-value error control, the
same sequential-testing machinery this project's \texttt{inquire}
workflow uses to classify evidence. POPPER terminates statistically, on
an error-rate bound over a population of tests; this work terminates
mechanically, on a deterministic kill predicate over a single binary
verdict, and persists the outcome as replayable memory where POPPER's
tests are ephemeral.

\section{Limitations}\label{limitations}

\emph{The mechanism evidence is existence-grade.} One audited
divergence, on one instance, in a program that was not preregistered
when it ran. The pilots' nulls are confounded by a selection artifact we
can name but not yet remove (the triage fast-path, \S\ref{null-regime}),
and the localization-hard band where the mechanism should live has not
been decisively tested; its one strong candidate did not reproduce at
HEAD. Nothing here is a rate.

\emph{The loop-versus-corpus control is unrun.} The externally-verified
arm hands the model both an execution loop and a labeled set of
must-flip cases; whether the lift comes from iterating against the gate
or from the labeled corpus alone is untested. The separating control, a
one-shot pass over the static cases with no loop, has not been run. (The
fair same-criteria battery is in place, and \#2501 was graded against
the committed probes at its own toolchain and matches the automated
fixes there; a full case-set co-grade of \#2501 at a matched toolchain
remains blocked by a toolchain mismatch.)

\emph{The smem is small and per-instance.} Hypothesis graphs in this
work are one markdown file per inquiry; cross-instance accumulation is
untested. That carries its own deflationary point: the file was never
the bottleneck at any repo size, so heavier stores need to earn their
keep at cross-instance scale, where the per-instance case never demanded
them.

\emph{How to refute this.} The central claims are built to fail loudly,
each against a committed artifact a hostile auditor runs rather than a
promise.

\begin{itemize}
\tightlist
\item
  \textbf{The mechanism claim} dies if the Verus receipt fails to
  replay, or if a discriminating program shows the externally-oracled
  fix wrong where the maintainer's general fix is right. The clean,
  forced-fresh dataset and both patches are committed to check, and the
  flux trail at \texttt{flux-1613-trail-v1} is there for the same test
  on the auditability case.
\item
  \textbf{The encoding-boundary claim}, that the oracle is not
  self-generable, dies if a self-graded model reaches the hard side with
  labels it authored independently of the predicate under test. That
  would show the oracle is inducible after all, and closing the gate's
  coverage on that side is the next experiment's preregistered target.
\item
  \textbf{The methodology claim} dies if, as audited cases accumulate,
  an external oracle stops separating from self-grading everywhere a
  receipt can see. The oracle's source is then not the active
  ingredient, and the mechanism README names that null as this thesis's
  own falsifier.
\item
  \textbf{The attribution} dies if the oracle bracket fails to replicate
  from its preregistered sample.
\item
  \textbf{The novelty claim}, a capability lift driven by a
  domain-general reasoning tool, dies if the lift is shown to ride on
  domain-specific tooling rather than the general operation, or if a
  documented prior result already demonstrated the same from a tool that
  carries no answer and no domain model (\S\ref{rw-scaffolds}). The
  comparative search (\S\ref{search}) publishes the queries to re-run.
\end{itemize}

\emph{Generator staleness.} The checkpoints are fixed (Sonnet 4.6 era),
and a more capable generator narrows the band where the mechanism is
observable, since it resolves more cases unaided. The mechanism claim
therefore survives only in the regime where verification is the
bottleneck, which is also the regime the discussion argues matters.

\section{Future work}\label{future-work}

The program reorganizes around the smem, in order of leverage.

\begin{itemize}
\tightlist
\item
  \textbf{The preregistered mechanism hunt.} The Verus and flux cases
  define the target band: localization-hard, verification-bottlenecked
  bugs where the symptom sits far from the cause and the suite cannot
  see the difference between fixes. Each loop now carries its deductive
  step before it runs: testing X, predict Y, refuted by Z
  (\texttt{METHODOLOGY-preregistration.md},
  \texttt{CANDIDATES-localization-hard.md}). Three to five audited
  existence cases, or the committed null, is the next paper.
\item
  \textbf{Coverage as the design lever, the oracle from approved
  history.} The Verus dissection (\S\ref{gate-general},
  \S\ref{enum-calib}) names two handles the next experiments turn
  directly. Gate \emph{coverage} sets the generalization frontier, so
  widening the enumeration to the genuine-divergence-preserve shape is a
  falsifiable test of whether the model then carves out the XOR's hard
  side. And the oracle a model cannot induce is mined from
  human-approved history, the merged fix, the regression suite, the
  resolved-issue label, so calibrating \texttt{comparator} differentially
  against base \emph{and} the approved fix is the concrete way to close
  a gate's blind spot without authoring a fresh oracle by hand.
  Prospectively, on a bug with no merged fix yet, the hard side has no
  correct verdict until a human spends the judgment.
\item
  \textbf{Conjecture: an agent cannot surprise itself.} Read the
  dissection one level up. Every warrant-producing step turns on a case
  the current hypothesis mishandles, the symmetric difference between
  what the inquirer believes and what is true. That case is the
  surprise, and the surprise is induction's to produce: it is what a
  test returns when the world disagrees with the deduced prediction. A
  test run against the agent's own belief returns no disagreement, so a
  self-grading agent generates no surprise, and abduction, which cannot
  fire without one, never re-fires. The agent supplies the leap and its
  deductions from within; the surprise it cannot, because induction
  against the self is not a test. That is why the externalized oracle
  was the active ingredient: it is the outside reference the inductive
  step needs and a single mind cannot be. Inquiry is the whole loop,
  abduction and deduction and induction together, and it completes only
  when its inductive step reaches past the agent; the external oracle is
  a constituent of that step, not a scaffold on it. The same pattern
  recurs across the companion posts (refutation in \emph{What Cannot Be
  False Cannot Be True}, stranger-replay in \emph{Verifiable Knowledge},
  the fold in \emph{Compress and Unfold}), which is suggestive rather
  than evidential and is kept here as a conjecture. We do not lean on it
  as a citation. It \emph{dies} if a recorded trace reaches the same
  held-out general predicate from self-authored labels only, with no
  external label, replay, human-approved reference, or environmental
  execution entering before the discriminating step; lucky guesses (no
  warrant) and already-solved cases (the difference already empty) do
  not count.
\item
  \textbf{Conjecture: in-context diff computation degrades with size.}
  Hold the harness fixed and the claim is about the operands: a model's
  ability to compute the XOR \emph{in its own context} should fall as
  the diff and case-space grow: a large structured operation attempted
  in prose rather than executed. Two regimes, one of them a conjecture.
  The \emph{endpoint} needs no experiment: once the operands exceed the
  context window, in-context computation is impossible, so
  externalization is the only option at scale, and a codebase outgrows
  any finite window eventually. The \emph{interior}, within the window
  but large, is the falsifiable part. Computing the XOR is a synthesis
  task, relating cases to one another to find the one the hypothesis
  mishandles, so the pairs to weigh grow with the square of the
  case-space and the load rises quadratically where a fact-lookup rises
  linearly; in-context accuracy should degrade as the operands grow, and
  faster than linearly. The underlying fall-off is the context-rot
  phenomenon (Chroma 2025); we have no direct measurement of XOR
  accuracy against diff size, and flag that rather than borrow a result
  that grades something else. It \emph{dies} if a within-harness sweep
  over diff and case-space size shows in-context XOR accuracy flat or
  rising with size. The Verus receipt is a single point on that axis.
\item
  \textbf{The hypothesis graph as a type.} Specify the abstract data
  structure independently of this harness: operations
  (perturb-and-classify, generate-edge-from-kill, prune, replay), the
  soundness invariant (every node reconstructible from its recorded
  trial), and a serialization any agent can emit. The goal is an
  interchange format for auditable reasoning, so that ``show your work''
  becomes a machine-checkable demand rather than a rhetorical one.
\item
  \textbf{Grade the trace, beyond the patch.} Bench design follows from
  the audit: report determinacy-aware denominators; build instruments
  whose oracle is expensive or absent and whose causes are hidden,
  because that is where method separates from reach; and score
  submissions on replayability, so a suite-green over-narrow fix
  (\S\ref{right-regime}) stops being indistinguishable from a root-cause
  one.
\item
  \textbf{Receipts-first contribution as the deployment lane.} The flux
  submission is the template: fix, discriminating receipt, soundness
  twin, full suite, replayable graph, residual flagged to the experts.
  Agent PRs are drowning in justified slop suspicion; the trace is the
  antidote, because it converts ``trust my patch'' into ``audit my
  ledger.'' Scaling that lane (and measuring maintainer response to
  trace-backed PRs) tests attestation-displacing-trust ecologically. The
  first maintainer merge of a trace-backed fix on a maintainer-stuck
  issue is the lane's golden ticket: an adversarial expert accepting the
  work on its ledger, with the doer invisible to the verdict.
\item
  \textbf{Cross-instance smem accumulation.} Let the graph grow across
  instances within a repo, then across repos within a domain; the
  current work tests the smem only at per-instance scope.
\item
  \textbf{Concurrent diagnosis over a shared graph.} The monotone graph
  (nodes append, kills idempotent) permits a lock-free fan-out: parallel
  agents on one shared graph re-verifying each other's kills instead of
  trusting them, with transitive accountability the precondition that
  keeps it safe, a conflict-free blackboard for the inquiry layered over
  git's blackboard for the code. Running it, and measuring the
  wall-clock win, is the to-do.
\end{itemize}

\section{Conclusion}\label{conclusion}

In the domain where every step is checkable, an agent's reasoning can be
made accountable rather than merely trusted. The hypothesis graph is a
replayable record that submits each generated step to a world-facing
trial and retains only what survives. The path from diagnosis to merge
becomes a sequence of falsifiable commitments a third party can rerun
rather than a verdict it must take on trust. This reduces the cost of
relying on an agent from reproducing its work to rerunning its record.
The guarantee holds only where the work is perturbable, where a trial
can be run and read; extending that region is left open, and the method
claims nothing beyond it.

Inside that checkable domain sits a bolder result. On a single
contamination-free bug, the externalized comparator carried Sonnet 4.6
to a fix Fable could not reach without it: a capability lift, not the
automation a scaffold usually buys. The within-model ablation attributes
it to the verdict source, the one input the model cannot author, with no
training involved. On a best-effort search we found no prior public
report of this specific crossing, a weaker model reaching a post-cutoff
fix a stronger one missed through a hidden external comparator, though
weaker-model-with-scaffold results exist on the automation axis.

This work is independently funded by the author. Every claim here ties
to a committed receipt, the nulls included.

\section{Availability and reproducibility}\label{availability}

\begin{itemize}
\tightlist
\item
  \textbf{Repositories.}
  \href{https://anonymous.4open.science/r/wepro-3D87}{anonymous.4open.science/r/wepro-3D87}
  (the bench run; frozen tags \texttt{prereg-pro-v1},
  \texttt{prereg-pro-v1-cheap}),
  \texttt{[redacted for double-blind review]}
  (prior-generation baseline, Zenodo-DOI'd),
  \href{https://anonymous.4open.science/r/weaudit-901E}{anonymous.4open.science/r/weaudit-901E}
  (the determinacy audit; every claim one row in \texttt{CLAIMS.md}, all
  728 verdicts in \texttt{COVERAGE.md}, mechanical spine re-derivable by
  grep),
  \texttt{[redacted for double-blind review]}
  (the audit as a portable tool for any SWE-bench-shaped bench;
  SWE-rebench run included),
  \href{https://anonymous.4open.science/r/hymech-117E}{anonymous.4open.science/r/hymech-117E}
  (the mechanism experiment; the Verus \#2219 lift with its clean,
  forced-fresh dataset and the 43,586-line climb trace, and the flux
  trail frozen at \texttt{flux-1613-trail-v1}; the README's artifact
  index maps every claim here to its committed path),
  \href{https://anonymous.4open.science/r/comparator-DA10}{anonymous.4open.science/r/comparator-DA10}
  (the externalized kill condition as a standalone, domain-general
  instrument: enumerate, calibrate against a known-good baseline, gate,
  with the \texttt{/debug} skill that drives the loop).
\item
  \textbf{Provenance artifacts.} Per-instance trajectories, hypothesis
  graphs, captured diffs, gate traces, and cost ledger under
  \texttt{runs/scored/artifacts/}; preregistrations at the freeze SHAs;
  the OSS ledger (\texttt{pr-receipts.jsonl}) with the GraphQL query
  that recomputes every number it asserts.
\item
  \textbf{Models and disclosure.} No training, fine-tuning, or learned
  weights in the harness; each stage drives a vendor's shipped agentic
  CLI over the same typed contracts. The bench run paired a Sonnet 4.6
  generator (extended thinking on) with a GPT-5.5 challenger (reasoning
  off, a point against reading the gate as model deliberation); the
  open-weight pair-swap ported the result wholesale, and each scored
  instance reads only its own artifacts. Versions and billing mode are
  in the cost ledger and the freeze SHAs.
\item
  \textbf{OSS deployment trace.} \textasciitilde385 hypothesis graphs at
  \texttt{[redacted for double-blind review]},
  one per investigated issue; PR-level outcomes pinned at
  \texttt{[redacted]}.
\item
  \textbf{Replication.} Boxes, budget, the per-instance cost ledger
  (\texttt{COST\_BASIS.md}), and the step-by-step rerun live in the run
  repo and the field guide \emph{How Not to Run SWE-bench Pro}.
\item
  \textbf{Companion writing.} The instrument story and field guide:
  \texttt{[redacted for double-blind review]}. The error corrected here, from the inside:
  \emph{Precisely Wrong}. The epistemology the discussion rests on:
  \emph{Verifiable Knowledge}
  (\texttt{[redacted for double-blind review]}) and its frame
  \emph{What Cannot Be False Cannot Be True}; its buildable-truth core
  in brief: \emph{Truth Is Buildable}. Dated provenance posts establish
  parallel rather than derivative development: \emph{Theory is
  load-bearing} (2026-03-17), \emph{The proof manual} (2026-04-05), and
  \emph{Type the question} (2026-04-08) predate ADI (2026-04-17);
  \emph{Evidence has a trajectory} (2026-04-27) and \emph{The Hypothesis
  Graph} (2026-04-28) predate CMM (2026-05-26).
\item
  \textbf{PDF.} Arxiv-shape build at
  \texttt{[redacted for double-blind review]},
  rebuilt from this markdown source by
  \texttt{[redacted for double-blind review]}; the source is
  canonical.
\item
  \textbf{DOI.} Prior-generation verified artifact:
  Zenodo-DOI'd.~Mechanism experiment
  (\href{https://anonymous.4open.science/r/hymech-117E}{hygraph-mechanism}):
  \texttt{[redacted for double-blind review]}.
  Pro bundle
  (\href{https://anonymous.4open.science/r/wepro-3D87}{swebench-pro}):
  \texttt{[redacted for double-blind review]}.
  This paper: DOI pending.
\item
  \textbf{License.} Skills released under \textbf{CC-BY-SA-NS}
  (\texttt{[redacted for double-blind review]});
  repo-level terms in each \texttt{LICENSE.md}. The harness an outsider
  clones is the same harness that produced the published numbers.
\end{itemize}

\textbf{Reproducibility invitation.} \emph{Nullius in verba.} Every
number here is recomputable from committed artifacts: the bench verdicts
by re-running the official grader on captured diffs, the audit's
mechanical spine by grep, the flux divergence by replaying the receipt
programs against both committed patches. Doubts should be filed as
issues against the relevant repository; confirmed corrections fold into
the next versioned artifact, as the retraction noted in
\S\ref{right-regime} and the reversal this version reports already have.

\section*{LLM collaboration
disclosure}\label{llm-collaboration-disclosure}
\addcontentsline{toc}{section}{LLM collaboration disclosure}

LLMs enter this work in three roles. \emph{Subject of study}: the
harness under evaluation uses frontier LLMs as generator and challenger,
with versions, billing mode, and provenance disclosed in
\S\ref{availability} and the artifacts. \emph{Instrument}: model pairs
adversarially verify the audit's two-expert tier (one constructs, an
independent family refutes) and blind-judge the mechanism pilots, always
with the mechanical layer (grep, grader, replay) holding the verdict.
\emph{Writing aid}: the prose was drafted and revised with Anthropic's
Claude (Opus 4.8 and Fable 5) from human-authored outlines and session
notes, with adversarial review from OpenAI's codex (GPT-5.5); the
claims, methodology, numbers, and argument structure are the author's.
No LLM decided what to publish.

\section*{Acknowledgments}\label{acknowledgments}
\addcontentsline{toc}{section}{Acknowledgments}

We thank John Laird for endorsing this submission, and the flux
maintainers for engaging with a stranger's receipts on their hardest
open issue.

\section*{References}\label{references}
\addcontentsline{toc}{section}{References}

Works cited above, consolidated. Entries for which the text or the
\emph{Extended intellectual lineage} gives only a short name or arXiv
identifier are listed at that detail rather than filled out, so nothing
here is reconstructed beyond what the paper states. The author's
companion essays are listed separately under \emph{Availability and
reproducibility} as lineage, not as entitlement.

\begin{itemize}
\tightlist
\item
  Abdaljalil et al.~(2025). \emph{Theorem-of-Thought}: typed
  abductive/deductive/inductive reasoning agents. arXiv:2506.07106.
\item
  Agarwal (2026). \emph{Refute-or-Promote}: adversarial defect
  discovery. arXiv:2604.19049.
\item
  Aleithan et al.~(2024). \emph{SWE-Bench+}: a manual audit of SWE-bench
  (solution leakage and weak tests).
\item
  Anokhin et al.~(2024/25). \emph{AriGraph}: knowledge-graph memory for
  LLM agents. arXiv:2407.04363.
\item
  Anthropic / Sumers, T., Yao, S., Narasimhan, K. \& Griffiths, T.
  (2023/24). ``Cognitive Architectures for Language Agents'' (CoALA).
  arXiv:2309.02427.
\item
  Bacon, F. (1620). \emph{Novum Organum}.
\item
  Bylander et al.~(1991). On the computational complexity of abduction.
  \emph{Artificial Intelligence} 49.
\item
  Calcagno, C., Distefano, D., O'Hearn, P. \& Yang, H. (2009).
  Compositional shape analysis by means of bi-abduction. \emph{POPL}.
  (Facebook Infer.)
\item
  Chroma (2025). Context rot: degradation of in-context reasoning with
  input length.
\item
  Clarke, E., Grumberg, O., Jha, S., Lu, Y. \& Veith, H. (2000).
  Counterexample-guided abstraction refinement (CEGAR). \emph{CAV}.
\item
  \emph{Confucius Code Agent} (2025). arXiv:2512.10398.
\item
  Cousot, P. \& Cousot, R. (1977). Abstract interpretation. \emph{POPL}.
\item
  de Kleer, J. (1986). ``An Assumption-based Truth Maintenance System.''
  \emph{Artificial Intelligence} 28.
\item
  de Kleer, J. \& Williams, B. (1987). Diagnosing multiple faults (the
  General Diagnostic Engine). \emph{Artificial Intelligence} 32.
\item
  Deng et al.~(2025). \emph{SWE-bench Pro}. arXiv:2509.16941.
\item
  Dewey, J. (1929). \emph{The Quest for Certainty}. Minton, Balch \& Co.
\item
  Doyle, J. (1979). ``A Truth Maintenance System.'' \emph{Artificial
  Intelligence} 12.
\item
  Dung, P. M. (1995). ``On the acceptability of arguments and its
  fundamental role in nonmonotonic reasoning, logic programming and
  n-person games.'' \emph{Artificial Intelligence} 77. (Abstract
  argumentation frameworks.)
\item
  Feynman, R. (1974). ``Cargo Cult Science.'' Caltech commencement
  address.
\item
  \emph{From Agent Traces to Trust} (2026): a provenance survey for
  agent reasoning. arXiv:2606.04990.
\item
  \emph{From Hypotheses to Factors} (2026): perturb-and-falsify over an
  append-only trace in quantitative finance. arXiv:2604.26747.
\item
  \emph{FVDebug} (2025): a hypothesis graph for debugging with
  model-selected nodes. arXiv:2510.15906.
\item
  Gilda, S. \& Gilda (2026). \emph{ADI}: a Peircean tripartite protocol
  over a symbolic knowledge graph. arXiv:2604.15727.
\item
  Goertzel, B., Iklé, M., Goertzel, I. \& Heljakka, A. (2008).
  \emph{Probabilistic Logic Networks}. Springer. (OpenCog / PLN.)
\item
  \emph{GradleFixer} (2025). arXiv:2510.08640.
\item
  Groth, P., Gibson, A. \& Velterop, J. (2010). ``The Anatomy of a
  Nanopublication.'' \emph{Information Services \& Use} 30.
\item
  He et al.~(2025). \emph{IDEA}: interactive rule learning with the
  three Peircean modes. ACL Findings. arXiv:2408.10455.
\item
  Jain et al.~(2024). \emph{LiveCodeBench}: temporal-holdout
  (post-cutoff) evaluation.
\item
  Jones, J., Harrold, M. J. \& Stasko, J. (2002). Spectrum-based fault
  localization. \emph{ICSE}.
\item
  Khalid \& Arora (2026). \emph{Cognitive Memory Manager} (CMM): a
  typed-node DAG mined from agent execution. OpenReview yCsHQnvvWY.
\item
  Kimball, A. W. (1957); Tukey, J. W. The Type III error: a precise
  answer to the wrong question.
\item
  Kirk, J., Wray, R. \& Laird, J. (2023). An LLM port of the Soar
  lineage. \emph{AAAI Spring Symposium}.
\item
  Laird, J., Newell, A. \& Rosenbloom, P. (1987). \emph{Soar}: an
  architecture for general intelligence. \emph{Artificial Intelligence}
  33.
\item
  Laurie, B., et al.~\emph{Certificate Transparency}. RFC 9162.
\item
  Le Goues, C., Nguyen, T., Forrest, S. \& Weimer, W. (2012). GenProg:
  search-based program repair. \emph{IEEE TSE}.
\item
  Liang et al.~(2023/24). \emph{Multi-Agent Debate}. arXiv:2305.19118.
\item
  Liblit, B., Naik, M., Zheng, A. X., Aiken, A. \& Jordan, M. I. (2005).
  Statistical debugging. \emph{PLDI}.
\item
  Liao et al.~(2026). \emph{BeliefMem}: candidate sets with Noisy-OR
  probabilistic update. arXiv:2605.05583.
\item
  Liu et al.~(2024). A survey of LLM-based agents for software
  engineering. arXiv:2409.02977.
\item
  Meehl, P. (1967). ``Theory-Testing in Psychology and Physics: A
  Methodological Paradox.'' \emph{Philosophy of Science} 34.
\item
  Modgil, S. \& Prakken, H. (2014). The ASPIC+ framework for structured
  argumentation. \emph{Argument \& Computation} 5.
\item
  Monperrus, M. (2018). ``Automatic Software Repair: A Bibliography.''
  \emph{ACM Computing Surveys} 51.
\item
  Moreau, L., et al.~(2013). \emph{PROV-DM}: the W3C provenance data
  model.
\item
  Necula, G. (1997). ``Proof-Carrying Code.'' \emph{POPL}.
\item
  O'Hearn, P. (2019). ``Incorrectness Logic.'' \emph{POPL} (separation
  logic).
\item
  \emph{OpenAI} (2026). SWE-bench Verified audit (February 2026): flawed
  tests and exact-gold-patch reproduction; recommends Pro.
\item
  \emph{ORACLE-SWE} (2026): ablating the oracle and specification
  signals that leak through a task. arXiv:2604.07789.
\item
  Pearl, J. (1988). \emph{Probabilistic Reasoning in Intelligent
  Systems}. Morgan Kaufmann. (Bayesian networks as DAGs.)
\item
  Pearl, J. (2000/2009). \emph{Causality: Models, Reasoning, and
  Inference}. Cambridge University Press.
\item
  Peirce, C. S. (1878). \emph{Illustrations of the Logic of Science}.
  \emph{Popular Science Monthly}. (The three-mode taxonomy of inquiry.)
\item
  Peirce, C. S. (1903). \emph{Pragmatism as the Logic of Abduction}
  (Harvard Lectures on Pragmatism).
\item
  \emph{POPPER} (2025): agentic sequential hypothesis testing under
  e-value control. arXiv:2502.09858.
\item
  Popper, K. (1934/1959). \emph{The Logic of Scientific Discovery}.
\item
  \emph{Portable Agent Memory} (2026): a Merkle-DAG provenance memory
  citing the Soar lineage. arXiv:2605.11032.
\item
  Ramsey, F. P. (1926). ``Truth and Probability.'' In \emph{The
  Foundations of Mathematics} (1931).
\item
  Reiter, R. (1987). ``A Theory of Diagnosis from First Principles.''
  \emph{Artificial Intelligence} 32.
\item
  Roychoudhury, A., et al.~(2025). Programming with trust: verification,
  testing, and analysis built into the agent. arXiv:2502.13767.
\item
  Solar-Lezama, A., et al.~(2006). Counterexample-guided inductive
  synthesis (CEGIS). \emph{ASPLOS}.
\item
  Stroebl, B., et al.~(2025). \emph{HAL}: a cost-transparent agent
  leaderboard.
\item
  \emph{SLUMP} (2026): an underspecified-by-design coding benchmark.
  arXiv:2603.17104.
\item
  \emph{SWE-Effi} (2025): effectiveness from scaffold-model synergy.
  arXiv:2509.09853.
\item
  \emph{SWE-rebench} (2025): post-cutoff filtering as a contamination
  strategy.
\item
  \emph{SWT-Bench} (2024): grading generated tests by fail-on-original
  then pass-after-golden-patch. arXiv:2406.12952.
\item
  \emph{Traxia} (2026): agent-native scientific publishing.
  arXiv:2606.08256.
\item
  Vovk, V. \& Wang, R. (2021). ``E-values: Calibration, combination, and
  applications.'' \emph{Annals of Statistics} 49.
\item
  Wald, A. (1947). \emph{Sequential Analysis}. Wiley.
\item
  Wang, P. (2013). \emph{Non-Axiomatic Logic: A Model of Intelligent
  Reasoning}. World Scientific. (NARS.)
\item
  Wang et al.~(2023). \emph{Voyager}: an open-ended embodied agent with
  a skill library. arXiv:2305.16291.
\item
  Wang et al.~(2024). \emph{OpenHands} (formerly OpenDevin): an open
  platform for AI software developers.
\item
  Wang et al.~(2025). A survey of agentic-programming systems
  (persistent structured memory as an open challenge). arXiv:2508.11126.
\item
  Wang, X., Pradel, M. \& Liu (2026). Plausible patches that pass tests
  yet diverge from developer intent. \emph{ICSE}.
\item
  Yang et al.~(2024). \emph{SWE-agent}: agent-computer interfaces for
  software engineering.
\item
  Yang et al.~(2026). \emph{CausaLab}: an evolving structural causal
  model in a DSL. arXiv:2605.26029.
\item
  Yao, S., et al.~(2023). \emph{ReAct}: synergizing reasoning and acting
  in language models.
\item
  Yehudai et al.~(2025). Trajectory-level assessment of agents.
  arXiv:2503.16416.
\item
  Zeller, A. \& Hildebrandt, R. (2002). ``Simplifying and Isolating
  Failure-Inducing Input'' (delta debugging). \emph{IEEE TSE} 28.
\item
  Zhang et al.~(2024/25). \emph{AutoCodeRover}: autonomous program
  improvement.
\item
  Zilberstein, N., Saliling \& Silva, A. (2024). Outcome Separation
  Logic; tri-abduction for branch composition. arXiv:2305.04842.
\end{itemize}

\section{Novelty and comparative search protocol}\label{search}

\begin{itemize}
\tightlist
\item
  \textbf{Why this section exists.} Several claims here take the form
  \emph{``we found no prior X.''} Such claims are only as honest as the
  search they rest on. This section publishes the queries so an auditor
  can re-run the search and either confirm the gap or find what we
  missed.
\item
  \textbf{Sources searched.} Google Scholar; arXiv (cs.LG, cs.AI, cs.CL,
  cs.SE); ACL Anthology; OpenReview; GitHub code and repository search;
  public SWE-bench leaderboard and submission archives.
\item
  \textbf{Queries by claim.}

  \begin{itemize}
  \tightlist
  \item
    \textbf{Peircean SE-agent loop.} \emph{``Peirce'' ``LLM agent''
    abduction deduction induction software engineering};
    \emph{``abduction deduction induction'' ``LLM agent'' ``software
    engineering''}; \emph{``Peirce'' ``SWE-bench'' agent}.
  \item
    \textbf{Hypothesis-graph agent memory.} \emph{``hypothesis graph''
    ``LLM agent'' memory}; \emph{``belief graph'' ``LLM agent'' memory
    hypothesis}; \emph{``knowledge graph'' ``LLM agent'' ``hypothesis''
    ``memory''}.
  \item
    \textbf{Blind multi-model hypothesis-stage filtering.}
    \emph{``multi-agent'' ``code'' ``hypothesis'' ``SWE-bench''};
    \emph{``blind'' ``multi-agent'' ``code review'' LLM}.
  \item
    \textbf{Trajectory-shape termination gates.} \emph{``LLM agent''
    termination criteria trajectory}; \emph{``finite state'' ``LLM
    agent'' ``termination''}.
  \item
    \textbf{Benchmark determinacy / construct-validity audits.}
    \emph{``SWE-bench Pro'' ``construct validity''};
    \emph{``underdetermined'' benchmark ``problem statement'' test};
    \emph{``specification'' ``ambiguity'' ``SWE-bench'' audit}.
  \item
    \textbf{Full per-instance provenance on SWE-bench Pro.}
    \emph{SWE-bench Pro leaderboard submissions trajectories cost
    ledger}; \emph{site:github.com ``swe-bench-pro'' ``trajs''}.
  \item
    \textbf{Sub-\$1k Pro replication and per-instance cost ledger.}
    \emph{``SWE-bench Pro'' ``cost per instance''};
    \emph{``SWE-rebench'' cost per problem}.
  \end{itemize}
\item
  \textbf{Caveats.} The search is best-effort and bounded by visible-web
  indexing; private industry work and non-indexed venues are not
  covered. Discoveries of overlapping prior work should be reported as
  issues for citation update.
\end{itemize}

\subsection{Comparative search supporting the artifact
claim}\label{comparative-search-supporting-the-artifact-claim}

The artifact claim (\S\ref{availability}), \emph{no method documented
publishes per-instance receipts at this depth on SWE-bench Pro},
requires a comparative search. The claim concerns receipts, with rate
set aside. The bar: published per-instance trajectories, captured diffs,
gate or evaluator traces, cost ledger, and reproducible run conditions.

\textbf{Candidate audit (against the receipt bar).} Each top public
submission or comparable report is checked for: published per-instance
trajectories (T), captured diffs (D), evaluator/gate traces (G),
per-instance cost ledger (C), reproducible frozen artifact (R), and
resolve rate at or above ours on the same bench. Receipt-bar columns are
\emph{present} (\cmark{}), \emph{partial} (\hvhalf{}), or \emph{absent}
(\xmark{}).

{\def\LTcaptype{none} % do not increment counter
\begin{longtable}[]{@{}
  >{\raggedright\arraybackslash}p{(\linewidth - 16\tabcolsep) * \real{0.2500}}
  >{\raggedright\arraybackslash}p{(\linewidth - 16\tabcolsep) * \real{0.1125}}
  >{\raggedright\arraybackslash}p{(\linewidth - 16\tabcolsep) * \real{0.0375}}
  >{\raggedright\arraybackslash}p{(\linewidth - 16\tabcolsep) * \real{0.0375}}
  >{\raggedright\arraybackslash}p{(\linewidth - 16\tabcolsep) * \real{0.0375}}
  >{\raggedright\arraybackslash}p{(\linewidth - 16\tabcolsep) * \real{0.0375}}
  >{\raggedright\arraybackslash}p{(\linewidth - 16\tabcolsep) * \real{0.0375}}
  >{\raggedright\arraybackslash}p{(\linewidth - 16\tabcolsep) * \real{0.2000}}
  >{\raggedright\arraybackslash}p{(\linewidth - 16\tabcolsep) * \real{0.2500}}@{}}
\toprule\noalign{}
\begin{minipage}[b]{\linewidth}\raggedright
Submission / report
\end{minipage} & \begin{minipage}[b]{\linewidth}\raggedright
Bench
\end{minipage} & \begin{minipage}[b]{\linewidth}\raggedright
T
\end{minipage} & \begin{minipage}[b]{\linewidth}\raggedright
D
\end{minipage} & \begin{minipage}[b]{\linewidth}\raggedright
G
\end{minipage} & \begin{minipage}[b]{\linewidth}\raggedright
C
\end{minipage} & \begin{minipage}[b]{\linewidth}\raggedright
R
\end{minipage} & \begin{minipage}[b]{\linewidth}\raggedright
Rate \(\geq\) ours
\end{minipage} & \begin{minipage}[b]{\linewidth}\raggedright
Notes
\end{minipage} \\
\midrule\noalign{}
\endhead
\bottomrule\noalign{}
\endlastfoot
Official \texttt{swebench/experiments} repo (multiple top entries) &
Verified & \cmark{} & \cmark{} & \xmark{} & \xmark{} & \hvhalf{} &
Various & Minimum publication norm: trajs/logs/patch.diff/report. No
gate traces, no cost ledger. \\
Top vendor leaderboard entries (Claude Code, OpenHands, SWE-agent,
AutoCodeRover) & Verified & \hvhalf{} & \hvhalf{} & \xmark{} & \xmark{}
& \xmark{} & Reported below 97\% & Submissions report numbers;
reproducible bundles and cost ledgers rarely published. \\
SWE-bench Pro official page (Scale) & Pro & \hvhalf{} & \hvhalf{} &
\xmark{} & \xmark{} & \xmark{} & N/A (curator) & Uncapped cost (250-turn
limit). No per-instance cost ledger. \\
Nilenso Pro trajectory analysis & Pro & \hvhalf{} & \xmark{} & \xmark{}
& \hvhalf{} & \xmark{} & N/A (third-party) & Cost/token/time analysis
across four frontier models. Not a submission. \\
SWE-rebench public reports & rebench & \hvhalf{} & \hvhalf{} & \xmark{}
& \cmark{} & \hvhalf{} & Below ours & Strong cost transparency (Cursor
Composer 2.5 at \$0.23/problem). \\
\textbf{This work: Verified} & Verified & \cmark{} & \cmark{} & \cmark{}
& \cmark{} & \cmark{} & 426 / 438 eligible (97.3\%) & Companion repo
\texttt{swebench-verified}; Zenodo DOI; gate traces and cost ledger
committed. \\
\textbf{This work: Pro} & Pro & \cmark{} & \cmark{} & \cmark{} &
\cmark{} & \cmark{} & 694/728 = 95.3\%; open-weight pair 678/728 =
93.1\% & Same frozen harness, whole eligible set, 0 incomplete; two
model pairs under one bundle. Public-split, gate-oracle regime: an
artifact claim, not a leaderboard claim (\S\ref{right-regime}). \\
\end{longtable}
}

\textbf{Reading.} No row above the two rows here combines all five
receipt-bar columns (T/D/G/C/R) on the same bench. The claim is about
receipt depth alone, leaving resolve rate out of it, and survives as
long as the table reads this way; a citation showing a fuller combined
receipt is the cleanest refutation.

\section{Extended intellectual lineage}\label{lineage}

\emph{Foundational sources grounding \S\ref{grounding},
\S\ref{epistemics}, \S\ref{hygraph}, and \S\ref{related-work}, collected
here so Related Work stays focused on contemporary systems.}

\subsection{Peircean inquiry and the philosophy of
science}\label{peircean-inquiry-and-the-philosophy-of-science}

\begin{itemize}
\tightlist
\item
  \textbf{Peirce 1878} (\emph{Illustrations of the Logic of Science}):
  the original three-mode taxonomy.
\item
  \textbf{Peirce 1903} (\emph{Pragmatism as the Logic of Abduction}):
  abduction as the only mode that introduces new content.
\item
  \textbf{Bacon 1620} (\emph{Novum Organum}); \textbf{Popper 1934}
  (\emph{The Logic of Scientific Discovery}): falsification as the
  inductive-side constraint.
\item
  \textbf{Ramsey 1926} (\emph{Truth and Probability}): operational
  credence as betting odds; the hypothesis graph's node-level semantics
  descends from this work.
\item
  \textbf{James 1907}; \textbf{Dewey 1929}: the pragmatist commitment
  that truth is inseparable from action.
\item
  \textbf{Meehl 1967}; \textbf{Feynman 1974} (``Cargo Cult Science''):
  the difference between rigor-shaped activity and actual rigor, the
  standard \S\ref{right-regime} holds itself to.
\item
  \textbf{Kimball 1957; Tukey}: the Type III error, the exact answer to
  the wrong question; the failure mode the determinacy audit
  (\S\ref{right-regime}) is built to prevent, examined in the companion
  post Precisely Wrong.
\end{itemize}

\subsection{The hypothesis graph's structural
ancestors}\label{the-hypothesis-graphs-structural-ancestors}

\begin{itemize}
\tightlist
\item
  \textbf{de Kleer 1986} (assumption-based truth-maintenance systems):
  persistent dependency structures over beliefs with mechanical
  retraction; the TMS keeps consistency where the hypothesis graph keeps
  \emph{trials}, and a TMS justification is not replayable by a
  stranger.
\item
  \textbf{Dung 1995} (abstract argumentation frameworks): attack
  relations between claims as first-class structure; the hypothesis
  graph's kill edges are attack edges bound to executions rather than
  arguments.
\item
  \textbf{Wald 1947} (sequential testing) and \textbf{Vovk \& Wang 2021}
  (e-values): the sequential-evidence framing that shaped
  \texttt{inquire}'s diagnostic stance. No accumulator is deployed: the
  code under test is deterministic, so the gate routes on the binary
  grader verdict (\S\ref{gating}).
\item
  \textbf{Pearl 1988; 2000/2009}: DAGs as the substrate for structured
  belief and causal-structure inference; this work borrows the
  typed-node/typed-edge form and leaves the probabilistic semantics
  behind.
\item
  \textbf{Zeller \& Hildebrandt 2002} (delta debugging): the canonical
  demonstration that mechanical perturbation of code is a productive
  inference primitive; optimization-shape where the hypothesis graph is
  methodology-shape.
\item
  \textbf{Reiter 1987} (\emph{A theory of diagnosis from first
  principles}) and \textbf{de Kleer \& Williams 1987} (the General
  Diagnostic Engine): model-based diagnosis as conflict sets, minimal
  hitting-set diagnoses, and entropy-minimizing choice of the next
  measurement. The hypothesis graph runs this abductive loop but over an
  open, LLM-generated hypothesis space rather than a fixed component
  model, and its ``where to perturb next'' is the GDE
  measurement-selection problem handed to the reasoner. Reiter's hitting
  sets are the theory of how multiple kills jointly localize a cause.
\item
  \textbf{Clarke et al.~2000} (counterexample-guided abstraction
  refinement, CEGAR), \textbf{Solar-Lezama et al.~2006}
  (counterexample-guided inductive synthesis, CEGIS), and \textbf{Cousot
  \& Cousot 1977} (abstract interpretation): the closest formal kin to
  the hypothesis graph's loop, in which a counterexample \emph{is} a
  kill that names the next refinement. The hypothesis graph generalizes
  that loop from a closed abstraction domain or synthesis grammar to an
  open hypothesis space, which forfeits the completeness and
  sound-by-construction refinement that closure buys (the open-world
  relevance boundary is undecidable in general, reducing to Rice's
  theorem) and recovers soundness only per node, through replay. Two
  pieces of their machinery are unported and natural to inherit: the
  real-versus-spurious counterexample check (a principled test of
  whether a kill is genuine or an artifact) and widening (a principled
  stop that over-approximates with stated precision loss, the formal
  form of declaring an inquiry's trajectory oscillatory and committing
  to the general shape).
\end{itemize}

\subsection{Bi-abductive and compositional
inference}\label{bi-abductive-and-compositional-inference}

\begin{itemize}
\tightlist
\item
  \textbf{Calcagno et al.~2009}: compositional shape analysis via
  bi-abduction; Facebook Infer as the industrial-scale instance of
  typed-mode inference on real code.
\item
  \textbf{O'Hearn 2019}: separation logic and incorrectness logic.
\item
  \textbf{Zilberstein, Saliling \& Silva 2024}
  (\href{https://arxiv.org/abs/2305.04842}{arXiv:2305.04842}): Outcome
  Separation Logic; tri-abduction for branch composition.
\item
  \textbf{Bylander et al.~1991}: abductive computational complexity.
\end{itemize}

\end{document}